% 2d_minimisation.tex
% Report on 2d minimisations algorithms (gradient descent and genetic algorithm).
% Vladimir Rutsky <altsysrq@gmail.com>
% 28.03.2009

\documentclass[10pt,a4paper,titlepage]{article}

% Spaces after commas.
\frenchspacing
% Minimal carrying number of characters,
\righthyphenmin=2

% Encoding support.
\usepackage{ucs}
\usepackage[utf8x]{inputenc}
\usepackage[T2A]{fontenc}
\usepackage[russian]{babel}

\usepackage{amsmath, amsthm, amssymb}

% From K.V.Voroncov Latex in samples, 2005.
\textheight=24cm   % Text height.
\textwidth=16cm    % Text width
\oddsidemargin=0pt % Left side indention.
\topmargin=-1.5cm  % Top side indention.
\parindent=24pt    % Paragraph indent.
\parskip=0pt       % Distance between paragraphs.
\tolerance=2000
%\flushbottom       % Page height aligning.

% Algorithms stuff.
\usepackage{algorithm}
\usepackage{algorithmic}

\newcommand{\INPUT}{\textbf{Вход:~}}
\newcommand{\OUTPUT}{\textbf{Выход:~}}
\floatname{algorithm}{Алгоритм}
\renewcommand{\algorithmicrepeat}{\textbf{do}}
\renewcommand{\algorithmicuntil}{\textbf{while}}

% Images.
%\usepackage[pdftex]{graphicx}

% Drawing figures.
\usepackage{subfig}

% Drawing graphs.
\usepackage{tikz}
\usetikzlibrary{trees,positioning,arrows}

% For listings.
\usepackage{listings}

% Transpose symbol in math mode.
\newcommand\T{^\mathrm{\textsf{\tiny{T}}}}
% Real number space.
\newcommand\RR{\mathbb{R}}

% Source code listings.
\renewcommand{\lstlistingname}{Исходный код}


\title{Решение задачи многомерной минимизации функции}
\author{Владимир Руцкий, 3057/2}
%\date{}

\begin{document}

\maketitle
\thispagestyle{empty}

% Content.
\section{Постановка задачи}
Требуется найти с наперёд заданной точностью точку в которой достигается минимум (локальный) многомерной функции $f(x)$ 
в некоторой области:
$$ \min f(x), \quad x \in \RR^n, $$
используя \textit{метод градиентного спуска} и \textit{генетический алгоритм}.

Исходная функция: $f(x) = x_1^3 + 2 x_2 + 4 \sqrt{2 + x_1^2 + x_2^2\mathstrut}
,$ заданная на $\RR^2.$ %TODO: Dimension specified directly.

\section{Исследование применимости методов}
\subsection{Метод градиентного спуска}
Исходная функция непрерывно дифференцируема:
$$
  \frac{\partial f}{\partial x_1} = 3 x_1^2 + \frac{4 x_1}{\sqrt{2 + x_1^2 + x_2^2\mathstrut}}
,
$$
$$
  \frac{\partial f}{\partial x_2} = ../../input/function_dx2.tex.
$$

Функция не является ни выпуклой, ни ограниченной на $\RR^2$, значит следует искать локальный минимум в некоторой области.

Построив график функции, можно попытаться найти достаточно близкую к локальному минимуму область, 
на которой функция будет выпуклой.
% \DeclareGraphicsRule{.pdftex}{pdf}{.pdftex}{}
% \includegraphics{data/function_graph.pdf}
\begin{figure}[h]
  \label{function-graph}
  %\caption{График функции $f(x) = x_1^3 + 2 x_2 + 4 \sqrt{2 + x_1^2 + x_2^2\mathstrut}
$}
  \caption{График функции $f(x)$}
  \begin{center}
  % GNUPLOT: LaTeX picture with Postscript
\begingroup%
\makeatletter%
\newcommand{\GNUPLOTspecial}{%
  \@sanitize\catcode`\%=14\relax\special}%
\setlength{\unitlength}{0.0500bp}%
\begin{picture}(5668,5668)(0,0)%
  {\GNUPLOTspecial{"
%!PS-Adobe-2.0 EPSF-2.0
%%Creator: gnuplot 4.2 patchlevel 4 
%%CreationDate: Tue Apr  7 03:24:19 2009
%%DocumentFonts: 
%%BoundingBox: 0 0 283 283
%%EndComments
%%BeginProlog
/gnudict 256 dict def
gnudict begin
%
% The following 6 true/false flags may be edited by hand if required
% The unit line width may also be changed
%
/Color false def
/Blacktext true def
/Solid false def
/Dashlength 1 def
/Landscape false def
/Level1 false def
/Rounded false def
/TransparentPatterns false def
/gnulinewidth 5.000 def
/userlinewidth gnulinewidth def
%
/vshift -66 def
/dl1 {
  10.0 Dashlength mul mul
  Rounded { currentlinewidth 0.75 mul sub dup 0 le { pop 0.01 } if } if
} def
/dl2 {
  10.0 Dashlength mul mul
  Rounded { currentlinewidth 0.75 mul add } if
} def
/hpt_ 31.5 def
/vpt_ 31.5 def
/hpt hpt_ def
/vpt vpt_ def
Level1 {} {
/SDict 10 dict def
systemdict /pdfmark known not {
  userdict /pdfmark systemdict /cleartomark get put
} if
SDict begin [
  /Title ()
  /Subject (gnuplot plot)
  /Creator (gnuplot 4.2 patchlevel 4 )
  /Author (Bob Rutsky,,,)
%  /Producer (gnuplot)
%  /Keywords ()
  /CreationDate (Tue Apr  7 03:24:19 2009)
  /DOCINFO pdfmark
end
} ifelse
%
% Gnuplot Prolog Version 4.2 (August 2006)
%
/M {moveto} bind def
/L {lineto} bind def
/R {rmoveto} bind def
/V {rlineto} bind def
/N {newpath moveto} bind def
/Z {closepath} bind def
/C {setrgbcolor} bind def
/f {rlineto fill} bind def
/vpt2 vpt 2 mul def
/hpt2 hpt 2 mul def
/Lshow {currentpoint stroke M 0 vshift R 
	Blacktext {gsave 0 setgray show grestore} {show} ifelse} def
/Rshow {currentpoint stroke M dup stringwidth pop neg vshift R
	Blacktext {gsave 0 setgray show grestore} {show} ifelse} def
/Cshow {currentpoint stroke M dup stringwidth pop -2 div vshift R 
	Blacktext {gsave 0 setgray show grestore} {show} ifelse} def
/UP {dup vpt_ mul /vpt exch def hpt_ mul /hpt exch def
  /hpt2 hpt 2 mul def /vpt2 vpt 2 mul def} def
/DL {Color {setrgbcolor Solid {pop []} if 0 setdash}
 {pop pop pop 0 setgray Solid {pop []} if 0 setdash} ifelse} def
/BL {stroke userlinewidth 2 mul setlinewidth
	Rounded {1 setlinejoin 1 setlinecap} if} def
/AL {stroke userlinewidth 2 div setlinewidth
	Rounded {1 setlinejoin 1 setlinecap} if} def
/UL {dup gnulinewidth mul /userlinewidth exch def
	dup 1 lt {pop 1} if 10 mul /udl exch def} def
/PL {stroke userlinewidth setlinewidth
	Rounded {1 setlinejoin 1 setlinecap} if} def
% Default Line colors
/LCw {1 1 1} def
/LCb {0 0 0} def
/LCa {0 0 0} def
/LC0 {1 0 0} def
/LC1 {0 1 0} def
/LC2 {0 0 1} def
/LC3 {1 0 1} def
/LC4 {0 1 1} def
/LC5 {1 1 0} def
/LC6 {0 0 0} def
/LC7 {1 0.3 0} def
/LC8 {0.5 0.5 0.5} def
% Default Line Types
/LTw {PL [] 1 setgray} def
/LTb {BL [] LCb DL} def
/LTa {AL [1 udl mul 2 udl mul] 0 setdash LCa setrgbcolor} def
/LT0 {PL [] LC0 DL} def
/LT1 {PL [4 dl1 2 dl2] LC1 DL} def
/LT2 {PL [2 dl1 3 dl2] LC2 DL} def
/LT3 {PL [1 dl1 1.5 dl2] LC3 DL} def
/LT4 {PL [6 dl1 2 dl2 1 dl1 2 dl2] LC4 DL} def
/LT5 {PL [3 dl1 3 dl2 1 dl1 3 dl2] LC5 DL} def
/LT6 {PL [2 dl1 2 dl2 2 dl1 6 dl2] LC6 DL} def
/LT7 {PL [1 dl1 2 dl2 6 dl1 2 dl2 1 dl1 2 dl2] LC7 DL} def
/LT8 {PL [2 dl1 2 dl2 2 dl1 2 dl2 2 dl1 2 dl2 2 dl1 4 dl2] LC8 DL} def
/Pnt {stroke [] 0 setdash gsave 1 setlinecap M 0 0 V stroke grestore} def
/Dia {stroke [] 0 setdash 2 copy vpt add M
  hpt neg vpt neg V hpt vpt neg V
  hpt vpt V hpt neg vpt V closepath stroke
  Pnt} def
/Pls {stroke [] 0 setdash vpt sub M 0 vpt2 V
  currentpoint stroke M
  hpt neg vpt neg R hpt2 0 V stroke
 } def
/Box {stroke [] 0 setdash 2 copy exch hpt sub exch vpt add M
  0 vpt2 neg V hpt2 0 V 0 vpt2 V
  hpt2 neg 0 V closepath stroke
  Pnt} def
/Crs {stroke [] 0 setdash exch hpt sub exch vpt add M
  hpt2 vpt2 neg V currentpoint stroke M
  hpt2 neg 0 R hpt2 vpt2 V stroke} def
/TriU {stroke [] 0 setdash 2 copy vpt 1.12 mul add M
  hpt neg vpt -1.62 mul V
  hpt 2 mul 0 V
  hpt neg vpt 1.62 mul V closepath stroke
  Pnt} def
/Star {2 copy Pls Crs} def
/BoxF {stroke [] 0 setdash exch hpt sub exch vpt add M
  0 vpt2 neg V hpt2 0 V 0 vpt2 V
  hpt2 neg 0 V closepath fill} def
/TriUF {stroke [] 0 setdash vpt 1.12 mul add M
  hpt neg vpt -1.62 mul V
  hpt 2 mul 0 V
  hpt neg vpt 1.62 mul V closepath fill} def
/TriD {stroke [] 0 setdash 2 copy vpt 1.12 mul sub M
  hpt neg vpt 1.62 mul V
  hpt 2 mul 0 V
  hpt neg vpt -1.62 mul V closepath stroke
  Pnt} def
/TriDF {stroke [] 0 setdash vpt 1.12 mul sub M
  hpt neg vpt 1.62 mul V
  hpt 2 mul 0 V
  hpt neg vpt -1.62 mul V closepath fill} def
/DiaF {stroke [] 0 setdash vpt add M
  hpt neg vpt neg V hpt vpt neg V
  hpt vpt V hpt neg vpt V closepath fill} def
/Pent {stroke [] 0 setdash 2 copy gsave
  translate 0 hpt M 4 {72 rotate 0 hpt L} repeat
  closepath stroke grestore Pnt} def
/PentF {stroke [] 0 setdash gsave
  translate 0 hpt M 4 {72 rotate 0 hpt L} repeat
  closepath fill grestore} def
/Circle {stroke [] 0 setdash 2 copy
  hpt 0 360 arc stroke Pnt} def
/CircleF {stroke [] 0 setdash hpt 0 360 arc fill} def
/C0 {BL [] 0 setdash 2 copy moveto vpt 90 450 arc} bind def
/C1 {BL [] 0 setdash 2 copy moveto
	2 copy vpt 0 90 arc closepath fill
	vpt 0 360 arc closepath} bind def
/C2 {BL [] 0 setdash 2 copy moveto
	2 copy vpt 90 180 arc closepath fill
	vpt 0 360 arc closepath} bind def
/C3 {BL [] 0 setdash 2 copy moveto
	2 copy vpt 0 180 arc closepath fill
	vpt 0 360 arc closepath} bind def
/C4 {BL [] 0 setdash 2 copy moveto
	2 copy vpt 180 270 arc closepath fill
	vpt 0 360 arc closepath} bind def
/C5 {BL [] 0 setdash 2 copy moveto
	2 copy vpt 0 90 arc
	2 copy moveto
	2 copy vpt 180 270 arc closepath fill
	vpt 0 360 arc} bind def
/C6 {BL [] 0 setdash 2 copy moveto
	2 copy vpt 90 270 arc closepath fill
	vpt 0 360 arc closepath} bind def
/C7 {BL [] 0 setdash 2 copy moveto
	2 copy vpt 0 270 arc closepath fill
	vpt 0 360 arc closepath} bind def
/C8 {BL [] 0 setdash 2 copy moveto
	2 copy vpt 270 360 arc closepath fill
	vpt 0 360 arc closepath} bind def
/C9 {BL [] 0 setdash 2 copy moveto
	2 copy vpt 270 450 arc closepath fill
	vpt 0 360 arc closepath} bind def
/C10 {BL [] 0 setdash 2 copy 2 copy moveto vpt 270 360 arc closepath fill
	2 copy moveto
	2 copy vpt 90 180 arc closepath fill
	vpt 0 360 arc closepath} bind def
/C11 {BL [] 0 setdash 2 copy moveto
	2 copy vpt 0 180 arc closepath fill
	2 copy moveto
	2 copy vpt 270 360 arc closepath fill
	vpt 0 360 arc closepath} bind def
/C12 {BL [] 0 setdash 2 copy moveto
	2 copy vpt 180 360 arc closepath fill
	vpt 0 360 arc closepath} bind def
/C13 {BL [] 0 setdash 2 copy moveto
	2 copy vpt 0 90 arc closepath fill
	2 copy moveto
	2 copy vpt 180 360 arc closepath fill
	vpt 0 360 arc closepath} bind def
/C14 {BL [] 0 setdash 2 copy moveto
	2 copy vpt 90 360 arc closepath fill
	vpt 0 360 arc} bind def
/C15 {BL [] 0 setdash 2 copy vpt 0 360 arc closepath fill
	vpt 0 360 arc closepath} bind def
/Rec {newpath 4 2 roll moveto 1 index 0 rlineto 0 exch rlineto
	neg 0 rlineto closepath} bind def
/Square {dup Rec} bind def
/Bsquare {vpt sub exch vpt sub exch vpt2 Square} bind def
/S0 {BL [] 0 setdash 2 copy moveto 0 vpt rlineto BL Bsquare} bind def
/S1 {BL [] 0 setdash 2 copy vpt Square fill Bsquare} bind def
/S2 {BL [] 0 setdash 2 copy exch vpt sub exch vpt Square fill Bsquare} bind def
/S3 {BL [] 0 setdash 2 copy exch vpt sub exch vpt2 vpt Rec fill Bsquare} bind def
/S4 {BL [] 0 setdash 2 copy exch vpt sub exch vpt sub vpt Square fill Bsquare} bind def
/S5 {BL [] 0 setdash 2 copy 2 copy vpt Square fill
	exch vpt sub exch vpt sub vpt Square fill Bsquare} bind def
/S6 {BL [] 0 setdash 2 copy exch vpt sub exch vpt sub vpt vpt2 Rec fill Bsquare} bind def
/S7 {BL [] 0 setdash 2 copy exch vpt sub exch vpt sub vpt vpt2 Rec fill
	2 copy vpt Square fill Bsquare} bind def
/S8 {BL [] 0 setdash 2 copy vpt sub vpt Square fill Bsquare} bind def
/S9 {BL [] 0 setdash 2 copy vpt sub vpt vpt2 Rec fill Bsquare} bind def
/S10 {BL [] 0 setdash 2 copy vpt sub vpt Square fill 2 copy exch vpt sub exch vpt Square fill
	Bsquare} bind def
/S11 {BL [] 0 setdash 2 copy vpt sub vpt Square fill 2 copy exch vpt sub exch vpt2 vpt Rec fill
	Bsquare} bind def
/S12 {BL [] 0 setdash 2 copy exch vpt sub exch vpt sub vpt2 vpt Rec fill Bsquare} bind def
/S13 {BL [] 0 setdash 2 copy exch vpt sub exch vpt sub vpt2 vpt Rec fill
	2 copy vpt Square fill Bsquare} bind def
/S14 {BL [] 0 setdash 2 copy exch vpt sub exch vpt sub vpt2 vpt Rec fill
	2 copy exch vpt sub exch vpt Square fill Bsquare} bind def
/S15 {BL [] 0 setdash 2 copy Bsquare fill Bsquare} bind def
/D0 {gsave translate 45 rotate 0 0 S0 stroke grestore} bind def
/D1 {gsave translate 45 rotate 0 0 S1 stroke grestore} bind def
/D2 {gsave translate 45 rotate 0 0 S2 stroke grestore} bind def
/D3 {gsave translate 45 rotate 0 0 S3 stroke grestore} bind def
/D4 {gsave translate 45 rotate 0 0 S4 stroke grestore} bind def
/D5 {gsave translate 45 rotate 0 0 S5 stroke grestore} bind def
/D6 {gsave translate 45 rotate 0 0 S6 stroke grestore} bind def
/D7 {gsave translate 45 rotate 0 0 S7 stroke grestore} bind def
/D8 {gsave translate 45 rotate 0 0 S8 stroke grestore} bind def
/D9 {gsave translate 45 rotate 0 0 S9 stroke grestore} bind def
/D10 {gsave translate 45 rotate 0 0 S10 stroke grestore} bind def
/D11 {gsave translate 45 rotate 0 0 S11 stroke grestore} bind def
/D12 {gsave translate 45 rotate 0 0 S12 stroke grestore} bind def
/D13 {gsave translate 45 rotate 0 0 S13 stroke grestore} bind def
/D14 {gsave translate 45 rotate 0 0 S14 stroke grestore} bind def
/D15 {gsave translate 45 rotate 0 0 S15 stroke grestore} bind def
/DiaE {stroke [] 0 setdash vpt add M
  hpt neg vpt neg V hpt vpt neg V
  hpt vpt V hpt neg vpt V closepath stroke} def
/BoxE {stroke [] 0 setdash exch hpt sub exch vpt add M
  0 vpt2 neg V hpt2 0 V 0 vpt2 V
  hpt2 neg 0 V closepath stroke} def
/TriUE {stroke [] 0 setdash vpt 1.12 mul add M
  hpt neg vpt -1.62 mul V
  hpt 2 mul 0 V
  hpt neg vpt 1.62 mul V closepath stroke} def
/TriDE {stroke [] 0 setdash vpt 1.12 mul sub M
  hpt neg vpt 1.62 mul V
  hpt 2 mul 0 V
  hpt neg vpt -1.62 mul V closepath stroke} def
/PentE {stroke [] 0 setdash gsave
  translate 0 hpt M 4 {72 rotate 0 hpt L} repeat
  closepath stroke grestore} def
/CircE {stroke [] 0 setdash 
  hpt 0 360 arc stroke} def
/Opaque {gsave closepath 1 setgray fill grestore 0 setgray closepath} def
/DiaW {stroke [] 0 setdash vpt add M
  hpt neg vpt neg V hpt vpt neg V
  hpt vpt V hpt neg vpt V Opaque stroke} def
/BoxW {stroke [] 0 setdash exch hpt sub exch vpt add M
  0 vpt2 neg V hpt2 0 V 0 vpt2 V
  hpt2 neg 0 V Opaque stroke} def
/TriUW {stroke [] 0 setdash vpt 1.12 mul add M
  hpt neg vpt -1.62 mul V
  hpt 2 mul 0 V
  hpt neg vpt 1.62 mul V Opaque stroke} def
/TriDW {stroke [] 0 setdash vpt 1.12 mul sub M
  hpt neg vpt 1.62 mul V
  hpt 2 mul 0 V
  hpt neg vpt -1.62 mul V Opaque stroke} def
/PentW {stroke [] 0 setdash gsave
  translate 0 hpt M 4 {72 rotate 0 hpt L} repeat
  Opaque stroke grestore} def
/CircW {stroke [] 0 setdash 
  hpt 0 360 arc Opaque stroke} def
/BoxFill {gsave Rec 1 setgray fill grestore} def
/Density {
  /Fillden exch def
  currentrgbcolor
  /ColB exch def /ColG exch def /ColR exch def
  /ColR ColR Fillden mul Fillden sub 1 add def
  /ColG ColG Fillden mul Fillden sub 1 add def
  /ColB ColB Fillden mul Fillden sub 1 add def
  ColR ColG ColB setrgbcolor} def
/BoxColFill {gsave Rec PolyFill} def
/PolyFill {gsave Density fill grestore grestore} def
/h {rlineto rlineto rlineto gsave closepath fill grestore} bind def
%
% PostScript Level 1 Pattern Fill routine for rectangles
% Usage: x y w h s a XX PatternFill
%	x,y = lower left corner of box to be filled
%	w,h = width and height of box
%	  a = angle in degrees between lines and x-axis
%	 XX = 0/1 for no/yes cross-hatch
%
/PatternFill {gsave /PFa [ 9 2 roll ] def
  PFa 0 get PFa 2 get 2 div add PFa 1 get PFa 3 get 2 div add translate
  PFa 2 get -2 div PFa 3 get -2 div PFa 2 get PFa 3 get Rec
  gsave 1 setgray fill grestore clip
  currentlinewidth 0.5 mul setlinewidth
  /PFs PFa 2 get dup mul PFa 3 get dup mul add sqrt def
  0 0 M PFa 5 get rotate PFs -2 div dup translate
  0 1 PFs PFa 4 get div 1 add floor cvi
	{PFa 4 get mul 0 M 0 PFs V} for
  0 PFa 6 get ne {
	0 1 PFs PFa 4 get div 1 add floor cvi
	{PFa 4 get mul 0 2 1 roll M PFs 0 V} for
 } if
  stroke grestore} def
%
/languagelevel where
 {pop languagelevel} {1} ifelse
 2 lt
	{/InterpretLevel1 true def}
	{/InterpretLevel1 Level1 def}
 ifelse
%
% PostScript level 2 pattern fill definitions
%
/Level2PatternFill {
/Tile8x8 {/PaintType 2 /PatternType 1 /TilingType 1 /BBox [0 0 8 8] /XStep 8 /YStep 8}
	bind def
/KeepColor {currentrgbcolor [/Pattern /DeviceRGB] setcolorspace} bind def
<< Tile8x8
 /PaintProc {0.5 setlinewidth pop 0 0 M 8 8 L 0 8 M 8 0 L stroke} 
>> matrix makepattern
/Pat1 exch def
<< Tile8x8
 /PaintProc {0.5 setlinewidth pop 0 0 M 8 8 L 0 8 M 8 0 L stroke
	0 4 M 4 8 L 8 4 L 4 0 L 0 4 L stroke}
>> matrix makepattern
/Pat2 exch def
<< Tile8x8
 /PaintProc {0.5 setlinewidth pop 0 0 M 0 8 L
	8 8 L 8 0 L 0 0 L fill}
>> matrix makepattern
/Pat3 exch def
<< Tile8x8
 /PaintProc {0.5 setlinewidth pop -4 8 M 8 -4 L
	0 12 M 12 0 L stroke}
>> matrix makepattern
/Pat4 exch def
<< Tile8x8
 /PaintProc {0.5 setlinewidth pop -4 0 M 8 12 L
	0 -4 M 12 8 L stroke}
>> matrix makepattern
/Pat5 exch def
<< Tile8x8
 /PaintProc {0.5 setlinewidth pop -2 8 M 4 -4 L
	0 12 M 8 -4 L 4 12 M 10 0 L stroke}
>> matrix makepattern
/Pat6 exch def
<< Tile8x8
 /PaintProc {0.5 setlinewidth pop -2 0 M 4 12 L
	0 -4 M 8 12 L 4 -4 M 10 8 L stroke}
>> matrix makepattern
/Pat7 exch def
<< Tile8x8
 /PaintProc {0.5 setlinewidth pop 8 -2 M -4 4 L
	12 0 M -4 8 L 12 4 M 0 10 L stroke}
>> matrix makepattern
/Pat8 exch def
<< Tile8x8
 /PaintProc {0.5 setlinewidth pop 0 -2 M 12 4 L
	-4 0 M 12 8 L -4 4 M 8 10 L stroke}
>> matrix makepattern
/Pat9 exch def
/Pattern1 {PatternBgnd KeepColor Pat1 setpattern} bind def
/Pattern2 {PatternBgnd KeepColor Pat2 setpattern} bind def
/Pattern3 {PatternBgnd KeepColor Pat3 setpattern} bind def
/Pattern4 {PatternBgnd KeepColor Landscape {Pat5} {Pat4} ifelse setpattern} bind def
/Pattern5 {PatternBgnd KeepColor Landscape {Pat4} {Pat5} ifelse setpattern} bind def
/Pattern6 {PatternBgnd KeepColor Landscape {Pat9} {Pat6} ifelse setpattern} bind def
/Pattern7 {PatternBgnd KeepColor Landscape {Pat8} {Pat7} ifelse setpattern} bind def
} def
%
%
%End of PostScript Level 2 code
%
/PatternBgnd {
  TransparentPatterns {} {gsave 1 setgray fill grestore} ifelse
} def
%
% Substitute for Level 2 pattern fill codes with
% grayscale if Level 2 support is not selected.
%
/Level1PatternFill {
/Pattern1 {0.250 Density} bind def
/Pattern2 {0.500 Density} bind def
/Pattern3 {0.750 Density} bind def
/Pattern4 {0.125 Density} bind def
/Pattern5 {0.375 Density} bind def
/Pattern6 {0.625 Density} bind def
/Pattern7 {0.875 Density} bind def
} def
%
% Now test for support of Level 2 code
%
Level1 {Level1PatternFill} {Level2PatternFill} ifelse
%
/Symbol-Oblique /Symbol findfont [1 0 .167 1 0 0] makefont
dup length dict begin {1 index /FID eq {pop pop} {def} ifelse} forall
currentdict end definefont pop
end
%%EndProlog
gnudict begin
gsave
0 0 translate
0.050 0.050 scale
0 setgray
newpath
gsave % colour palette begin
/maxcolors 0 def
/HSV2RGB {  exch dup 0.0 eq {pop exch pop dup dup} % achromatic gray
  { /HSVs exch def /HSVv exch def 6.0 mul dup floor dup 3 1 roll sub
     /HSVf exch def /HSVi exch cvi def /HSVp HSVv 1.0 HSVs sub mul def
	 /HSVq HSVv 1.0 HSVs HSVf mul sub mul def 
	 /HSVt HSVv 1.0 HSVs 1.0 HSVf sub mul sub mul def
	 /HSVi HSVi 6 mod def 0 HSVi eq {HSVv HSVt HSVp}
	 {1 HSVi eq {HSVq HSVv HSVp}{2 HSVi eq {HSVp HSVv HSVt}
	 {3 HSVi eq {HSVp HSVq HSVv}{4 HSVi eq {HSVt HSVp HSVv}
	 {HSVv HSVp HSVq} ifelse} ifelse} ifelse} ifelse} ifelse
  } ifelse} def
/Constrain {
  dup 0 lt {0 exch pop}{dup 1 gt {1 exch pop} if} ifelse} def
/YIQ2RGB {
  3 copy -1.702 mul exch -1.105 mul add add Constrain 4 1 roll
  3 copy -0.647 mul exch -0.272 mul add add Constrain 5 1 roll
  0.621 mul exch -0.956 mul add add Constrain 3 1 roll } def
/CMY2RGB {  1 exch sub exch 1 exch sub 3 2 roll 1 exch sub 3 1 roll exch } def
/XYZ2RGB {  3 copy -0.9017 mul exch -0.1187 mul add exch 0.0585 mul exch add
  Constrain 4 1 roll 3 copy -0.0279 mul exch 1.999 mul add exch
  -0.9844 mul add Constrain 5 1 roll -0.2891 mul exch -0.5338 mul add
  exch 1.91 mul exch add Constrain 3 1 roll} def
/SelectSpace {ColorSpace (HSV) eq {HSV2RGB}{ColorSpace (XYZ) eq {
  XYZ2RGB}{ColorSpace (CMY) eq {CMY2RGB}{ColorSpace (YIQ) eq {YIQ2RGB}
  if} ifelse} ifelse} ifelse} def
/pm3dround {maxcolors 0 gt {dup 1 ge
	{pop 1} {maxcolors mul floor maxcolors 1 sub div} ifelse} if} def
/pm3dGamma 1.0 1.5 div def
/ColorSpace (RGB) def
Color false and { % COLOUR vs. GRAY map
  InterpolatedColor { %% Interpolation vs. RGB-Formula
    /g {stroke pm3dround /grayv exch def interpolate
        SelectSpace setrgbcolor} bind def
  }{
  /g {stroke pm3dround dup cF7 Constrain exch dup cF5 Constrain exch cF15 Constrain 
       SelectSpace setrgbcolor} bind def
  } ifelse
}{
  /g {stroke pm3dround pm3dGamma exp setgray} bind def
} ifelse
1.000 UL
LTb
1.000 UL
LTb
1119 2892 M
2828 448 V
4549 1228 M
3947 3340 L
1119 2892 M
0 1748 V
3947 3340 M
0 1749 V
4549 1228 M
0 1748 V
1119 4640 M
2828 449 V
4549 2976 M
3947 5089 L
1721 779 M
-14 47 V
1119 2892 M
14 -48 V
2192 853 M
-13 48 V
1591 2966 M
14 -47 V
2664 928 M
-14 48 V
2062 3041 M
14 -48 V
3134 1003 M
-13 48 V
2534 3115 M
13 -48 V
3606 1078 M
-14 48 V
3004 3190 M
14 -48 V
4077 1153 M
-14 48 V
3476 3265 M
13 -48 V
4549 1228 M
-14 48 V
3947 3340 M
14 -48 V
1721 779 M
61 9 V
2767 440 R
-62 -10 V
1621 1131 M
61 10 V
2766 439 R
-61 -10 V
1520 1483 M
62 10 V
2766 439 R
-62 -10 V
1420 1835 M
62 10 V
2766 439 R
-62 -9 V
1320 2187 M
62 10 V
2766 440 R
-62 -10 V
1220 2540 M
61 9 V
2766 439 R
-61 -10 V
1119 2892 M
62 10 V
2766 438 R
-61 -10 V
1119 2892 M
63 0 V
-63 249 R
63 0 V
-63 250 R
63 0 V
-63 250 R
63 0 V
-63 249 R
63 0 V
-63 250 R
63 0 V
-63 250 R
63 0 V
-63 250 R
63 0 V
stroke
LCb setrgbcolor
LTb
1.000 UP
1.000 UL
LT0
0.00 0.00 0.00 C 0.00 0.00 0.00 C 1720 922 M
118 160 V
117 135 V
118 112 V
118 92 V
118 75 V
118 59 V
118 47 V
118 36 V
118 28 V
117 23 V
118 19 V
118 18 V
118 20 V
118 24 V
118 29 V
118 39 V
117 49 V
118 62 V
118 77 V
118 95 V
118 114 V
118 137 V
118 161 V
118 188 V
1695 1004 M
118 159 V
117 134 V
118 112 V
118 92 V
118 74 V
118 59 V
118 46 V
118 36 V
117 28 V
118 22 V
118 19 V
118 18 V
118 20 V
118 24 V
118 30 V
118 38 V
117 50 V
118 62 V
118 78 V
118 95 V
118 115 V
118 137 V
118 162 V
117 188 V
1670 1086 M
117 158 V
118 134 V
118 111 V
118 92 V
118 73 V
118 59 V
118 45 V
118 36 V
117 27 V
118 22 V
118 19 V
118 19 V
118 19 V
118 24 V
118 31 V
118 38 V
117 50 V
118 63 V
118 78 V
118 96 V
118 116 V
118 137 V
118 163 V
117 189 V
1645 1169 M
117 158 V
118 133 V
118 110 V
118 91 V
118 73 V
118 58 V
118 45 V
118 35 V
117 27 V
118 22 V
118 19 V
118 18 V
118 20 V
118 24 V
118 31 V
117 39 V
118 50 V
118 63 V
118 79 V
118 96 V
118 117 V
118 138 V
118 163 V
117 189 V
1620 1252 M
117 157 V
118 133 V
118 110 V
118 90 V
stroke 2091 1742 M
118 72 V
118 57 V
118 45 V
117 35 V
118 26 V
118 22 V
118 19 V
118 18 V
118 20 V
118 24 V
118 31 V
117 39 V
118 51 V
118 64 V
118 79 V
118 97 V
118 117 V
118 139 V
117 164 V
118 190 V
1594 1336 M
118 157 V
118 132 V
118 109 V
118 89 V
118 72 V
118 57 V
118 44 V
117 34 V
118 26 V
118 21 V
118 19 V
118 18 V
118 20 V
118 24 V
117 32 V
118 40 V
118 51 V
118 64 V
118 80 V
118 98 V
118 117 V
118 140 V
117 164 V
118 191 V
1569 1421 M
118 156 V
118 131 V
118 109 V
118 89 V
118 71 V
118 56 V
117 43 V
118 34 V
118 25 V
118 21 V
118 18 V
118 18 V
118 21 V
118 24 V
117 32 V
118 40 V
118 52 V
118 65 V
118 81 V
118 98 V
118 118 V
117 141 V
118 165 V
118 191 V
1544 1507 M
118 156 V
118 130 V
118 108 V
118 88 V
118 70 V
118 55 V
117 43 V
118 33 V
118 25 V
118 21 V
118 18 V
118 18 V
118 20 V
118 25 V
117 32 V
118 41 V
118 52 V
118 66 V
118 81 V
118 99 V
118 119 V
117 141 V
118 166 V
118 192 V
1519 1594 M
118 155 V
118 130 V
118 107 V
118 87 V
118 69 V
118 55 V
117 42 V
118 32 V
stroke 2462 2271 M
118 25 V
118 20 V
118 17 V
118 18 V
118 21 V
117 25 V
118 32 V
118 42 V
118 53 V
118 66 V
118 82 V
118 100 V
118 120 V
117 142 V
118 166 V
118 193 V
1494 1682 M
118 155 V
118 129 V
118 106 V
118 86 V
118 69 V
117 54 V
118 41 V
118 31 V
118 24 V
118 20 V
118 17 V
118 18 V
118 20 V
117 26 V
118 33 V
118 42 V
118 54 V
118 67 V
118 83 V
118 101 V
117 120 V
118 143 V
118 167 V
118 193 V
1469 1772 M
118 153 V
118 129 V
118 105 V
118 86 V
118 68 V
117 53 V
118 40 V
118 31 V
118 23 V
118 19 V
118 17 V
118 18 V
117 20 V
118 26 V
118 34 V
118 42 V
118 55 V
118 68 V
118 84 V
118 101 V
117 122 V
118 143 V
118 168 V
118 194 V
1444 1863 M
118 153 V
118 128 V
118 105 V
118 84 V
117 67 V
118 52 V
118 40 V
118 30 V
118 22 V
118 19 V
118 16 V
118 18 V
117 21 V
118 26 V
118 34 V
118 43 V
118 56 V
118 69 V
118 84 V
117 103 V
118 122 V
118 144 V
118 168 V
118 195 V
1419 1956 M
118 152 V
118 127 V
118 104 V
118 84 V
117 67 V
118 51 V
118 39 V
118 29 V
118 22 V
118 17 V
118 16 V
118 18 V
stroke 2834 2682 M
117 21 V
118 26 V
118 35 V
118 44 V
118 57 V
118 69 V
118 86 V
117 103 V
118 123 V
118 144 V
118 169 V
118 195 V
1394 2050 M
118 152 V
118 126 V
118 104 V
118 84 V
117 65 V
118 51 V
118 38 V
118 28 V
118 21 V
118 17 V
118 16 V
117 17 V
118 21 V
118 27 V
118 36 V
118 45 V
118 57 V
118 70 V
118 86 V
117 104 V
118 123 V
118 146 V
118 169 V
118 196 V
1369 2146 M
118 151 V
118 127 V
118 103 V
117 83 V
118 65 V
118 50 V
118 38 V
118 27 V
118 21 V
118 16 V
118 16 V
117 17 V
118 21 V
118 27 V
118 36 V
118 46 V
118 57 V
118 71 V
117 87 V
118 104 V
118 124 V
118 146 V
118 170 V
118 195 V
1344 2243 M
118 152 V
118 126 V
118 103 V
117 83 V
118 65 V
118 49 V
118 37 V
118 27 V
118 21 V
118 16 V
117 15 V
118 17 V
118 21 V
118 28 V
118 36 V
118 46 V
118 58 V
118 72 V
117 87 V
118 104 V
118 124 V
118 146 V
118 170 V
118 196 V
1319 2343 M
118 151 V
118 126 V
117 103 V
118 83 V
118 65 V
118 49 V
118 37 V
118 27 V
118 20 V
118 16 V
117 15 V
118 17 V
118 21 V
118 28 V
118 37 V
118 46 V
stroke 3205 3184 M
118 58 V
117 72 V
118 87 V
118 104 V
118 125 V
118 146 V
118 170 V
118 196 V
1294 2444 M
118 152 V
118 126 V
117 103 V
118 83 V
118 65 V
118 49 V
118 37 V
118 27 V
118 21 V
117 16 V
118 15 V
118 17 V
118 21 V
118 28 V
118 36 V
118 46 V
118 58 V
117 72 V
118 87 V
118 104 V
118 124 V
118 146 V
118 170 V
118 196 V
1269 2548 M
118 151 V
118 127 V
117 103 V
118 83 V
118 65 V
118 50 V
118 38 V
118 27 V
118 21 V
117 16 V
118 16 V
118 17 V
118 21 V
118 28 V
118 35 V
118 46 V
118 58 V
117 71 V
118 86 V
118 105 V
118 123 V
118 146 V
118 170 V
118 196 V
1244 2653 M
118 152 V
117 127 V
118 103 V
118 84 V
118 66 V
118 50 V
118 38 V
118 28 V
118 21 V
117 18 V
118 15 V
118 17 V
118 22 V
118 27 V
118 35 V
118 45 V
117 57 V
118 70 V
118 87 V
118 103 V
118 124 V
118 145 V
118 169 V
118 196 V
1219 2760 M
118 152 V
117 127 V
118 105 V
118 84 V
118 66 V
118 51 V
118 39 V
118 29 V
117 22 V
118 18 V
118 16 V
118 17 V
118 21 V
118 27 V
118 34 V
118 45 V
117 56 V
118 70 V
118 85 V
118 103 V
stroke 3576 3927 M
118 123 V
118 145 V
118 169 V
117 195 V
1194 2868 M
117 153 V
118 128 V
118 105 V
118 85 V
118 67 V
118 52 V
118 39 V
118 30 V
117 23 V
118 18 V
118 17 V
118 17 V
118 21 V
118 26 V
118 34 V
117 44 V
118 55 V
118 69 V
118 85 V
118 102 V
118 122 V
118 144 V
118 169 V
117 194 V
1169 2979 M
117 153 V
118 128 V
118 106 V
118 86 V
118 67 V
118 53 V
118 41 V
117 30 V
118 24 V
118 19 V
118 17 V
118 17 V
118 21 V
118 26 V
118 33 V
117 43 V
118 55 V
118 68 V
118 83 V
118 102 V
118 121 V
118 144 V
118 167 V
117 194 V
1144 3090 M
117 154 V
118 129 V
118 107 V
118 86 V
118 69 V
118 53 V
118 42 V
117 31 V
118 24 V
118 20 V
118 17 V
118 18 V
118 20 V
118 26 V
118 32 V
117 43 V
118 53 V
118 68 V
118 83 V
118 100 V
118 121 V
118 143 V
117 167 V
118 193 V
1118 3203 M
118 155 V
118 129 V
118 108 V
118 87 V
118 69 V
118 55 V
118 42 V
117 32 V
118 25 V
118 20 V
118 17 V
118 18 V
118 20 V
118 26 V
117 32 V
118 42 V
118 53 V
118 66 V
118 82 V
118 100 V
118 120 V
118 142 V
117 166 V
118 193 V
stroke 3947 5002 M
1720 922 M
-25 82 V
-25 82 V
-25 83 V
-25 83 V
-26 84 V
-25 85 V
-25 86 V
-25 87 V
-25 88 V
-25 90 V
-25 91 V
-25 93 V
-25 94 V
-25 96 V
-25 97 V
-25 100 V
-25 101 V
-25 104 V
-25 105 V
-25 107 V
-25 108 V
-25 111 V
-25 111 V
-26 113 V
1838 1082 M
-25 81 V
-26 81 V
-25 83 V
-25 82 V
-25 84 V
-25 84 V
-25 86 V
-25 86 V
-25 88 V
-25 88 V
-25 91 V
-25 92 V
-25 94 V
-25 95 V
-25 98 V
-25 99 V
-25 102 V
-25 103 V
-25 106 V
-25 107 V
-26 109 V
-25 111 V
-25 112 V
-25 114 V
1955 1217 M
-25 80 V
-25 81 V
-25 82 V
-25 82 V
-25 83 V
-25 83 V
-25 85 V
-25 86 V
-25 87 V
-25 88 V
-25 90 V
-25 91 V
-25 93 V
-25 96 V
-25 97 V
-25 99 V
-25 102 V
-25 104 V
-26 106 V
-25 107 V
-25 110 V
-25 111 V
-25 113 V
-25 114 V
2073 1329 M
-25 80 V
-25 80 V
-25 81 V
-25 82 V
-25 82 V
-25 83 V
-25 84 V
-25 85 V
-25 86 V
-25 87 V
-25 90 V
-25 90 V
-25 93 V
-25 95 V
-25 97 V
-26 99 V
-25 102 V
-25 104 V
-25 106 V
-25 109 V
-25 110 V
-25 112 V
-25 114 V
-25 115 V
2191 1421 M
-25 80 V
-25 80 V
-25 80 V
stroke 2116 1661 M
-25 81 V
-25 81 V
-25 83 V
-25 83 V
-25 84 V
-25 85 V
-25 87 V
-25 88 V
-25 90 V
-25 93 V
-26 94 V
-25 97 V
-25 99 V
-25 102 V
-25 104 V
-25 107 V
-25 109 V
-25 111 V
-25 113 V
-25 114 V
-25 116 V
2309 1496 M
-25 79 V
-25 79 V
-25 80 V
-25 80 V
-25 81 V
-25 82 V
-25 82 V
-25 83 V
-25 85 V
-25 86 V
-26 87 V
-25 90 V
-25 91 V
-25 94 V
-25 97 V
-25 99 V
-25 102 V
-25 104 V
-25 108 V
-25 109 V
-25 112 V
-25 113 V
-25 116 V
-25 116 V
2427 1555 M
-25 79 V
-25 79 V
-25 79 V
-25 79 V
-25 81 V
-25 81 V
-25 81 V
-25 83 V
-26 84 V
-25 85 V
-25 86 V
-25 89 V
-25 91 V
-25 93 V
-25 96 V
-25 99 V
-25 102 V
-25 105 V
-25 108 V
-25 110 V
-25 113 V
-25 114 V
-25 116 V
-25 118 V
2545 1602 M
-25 78 V
-25 78 V
-25 79 V
-25 79 V
-25 80 V
-26 80 V
-25 81 V
-25 82 V
-25 83 V
-25 84 V
-25 86 V
-25 88 V
-25 90 V
-25 93 V
-25 95 V
-25 99 V
-25 102 V
-25 106 V
-25 108 V
-25 111 V
-25 113 V
-25 116 V
-25 117 V
-25 118 V
2663 1638 M
-25 78 V
-25 78 V
-25 78 V
-26 79 V
-25 79 V
-25 80 V
-25 80 V
stroke 2487 2190 M
-25 81 V
-25 82 V
-25 84 V
-25 85 V
-25 87 V
-25 89 V
-25 92 V
-25 95 V
-25 99 V
-25 102 V
-25 106 V
-25 109 V
-25 112 V
-25 114 V
-26 116 V
-25 118 V
-25 119 V
2781 1666 M
-26 78 V
-25 77 V
-25 78 V
-25 78 V
-25 79 V
-25 79 V
-25 80 V
-25 81 V
-25 81 V
-25 83 V
-25 84 V
-25 87 V
-25 88 V
-25 92 V
-25 95 V
-25 98 V
-25 103 V
-25 106 V
-25 109 V
-26 113 V
-25 115 V
-25 117 V
-25 118 V
-25 120 V
2898 1689 M
-25 77 V
-25 77 V
-25 78 V
-25 78 V
-25 78 V
-25 79 V
-25 80 V
-25 80 V
-25 81 V
-25 82 V
-25 84 V
-25 85 V
-25 88 V
-25 91 V
-25 95 V
-25 98 V
-26 103 V
-25 106 V
-25 111 V
-25 113 V
-25 115 V
-25 118 V
-25 119 V
-25 120 V
3016 1708 M
-25 77 V
-25 77 V
-25 78 V
-25 78 V
-25 78 V
-25 78 V
-25 80 V
-25 79 V
-25 81 V
-25 82 V
-25 83 V
-25 85 V
-25 88 V
-25 91 V
-26 94 V
-25 98 V
-25 103 V
-25 107 V
-25 110 V
-25 114 V
-25 116 V
-25 118 V
-25 119 V
-25 120 V
3134 1726 M
-25 77 V
-25 78 V
-25 77 V
-25 78 V
-25 78 V
-25 78 V
-25 80 V
-25 79 V
-25 81 V
-25 82 V
-25 83 V
stroke 2859 2597 M
-25 85 V
-26 87 V
-25 91 V
-25 94 V
-25 98 V
-25 103 V
-25 107 V
-25 110 V
-25 114 V
-25 116 V
-25 118 V
-25 120 V
-25 120 V
3252 1746 M
-25 77 V
-25 77 V
-25 78 V
-25 78 V
-25 78 V
-25 79 V
-25 79 V
-25 80 V
-25 80 V
-26 82 V
-25 84 V
-25 85 V
-25 87 V
-25 91 V
-25 94 V
-25 98 V
-25 103 V
-25 107 V
-25 111 V
-25 113 V
-25 116 V
-25 118 V
-25 119 V
-25 120 V
3370 1770 M
-25 77 V
-25 77 V
-25 78 V
-25 78 V
-25 78 V
-25 79 V
-25 80 V
-26 80 V
-25 81 V
-25 82 V
-25 84 V
-25 85 V
-25 88 V
-25 91 V
-25 95 V
-25 98 V
-25 103 V
-25 107 V
-25 110 V
-25 113 V
-25 115 V
-25 118 V
-25 119 V
-25 120 V
3488 1799 M
-25 78 V
-25 78 V
-25 78 V
-25 78 V
-26 79 V
-25 79 V
-25 80 V
-25 80 V
-25 82 V
-25 83 V
-25 84 V
-25 86 V
-25 89 V
-25 91 V
-25 95 V
-25 99 V
-25 102 V
-25 106 V
-25 110 V
-25 112 V
-25 115 V
-25 117 V
-25 118 V
-26 120 V
3606 1838 M
-25 77 V
-25 78 V
-26 79 V
-25 78 V
-25 80 V
-25 79 V
-25 81 V
-25 81 V
-25 82 V
-25 83 V
-25 85 V
-25 87 V
-25 90 V
-25 92 V
-25 95 V
stroke 3230 3085 M
-25 99 V
-25 102 V
-25 106 V
-25 109 V
-25 112 V
-26 114 V
-25 116 V
-25 118 V
-25 119 V
3723 1887 M
-25 78 V
-25 78 V
-25 79 V
-25 79 V
-25 80 V
-25 80 V
-25 81 V
-25 82 V
-25 83 V
-25 84 V
-25 86 V
-25 88 V
-25 90 V
-25 92 V
-25 96 V
-25 99 V
-25 102 V
-25 106 V
-26 108 V
-25 111 V
-25 113 V
-25 116 V
-25 116 V
-25 119 V
3841 1949 M
-25 78 V
-25 79 V
-25 79 V
-25 80 V
-25 80 V
-25 81 V
-25 82 V
-25 82 V
-25 84 V
-25 85 V
-25 87 V
-25 88 V
-25 91 V
-25 93 V
-25 97 V
-26 99 V
-25 102 V
-25 105 V
-25 107 V
-25 111 V
-25 112 V
-25 115 V
-25 116 V
-25 117 V
3959 2026 M
-25 79 V
-25 79 V
-25 80 V
-25 80 V
-25 81 V
-25 82 V
-25 82 V
-25 83 V
-25 85 V
-25 86 V
-25 87 V
-25 90 V
-25 91 V
-26 94 V
-25 97 V
-25 99 V
-25 102 V
-25 104 V
-25 108 V
-25 109 V
-25 112 V
-25 113 V
-25 116 V
-25 116 V
4077 2121 M
-25 79 V
-25 80 V
-25 80 V
-25 81 V
-25 82 V
-25 82 V
-25 83 V
-25 84 V
-25 86 V
-25 86 V
-26 89 V
-25 90 V
-25 92 V
-25 94 V
-25 97 V
-25 99 V
-25 102 V
-25 105 V
-25 106 V
stroke 3601 3818 M
-25 109 V
-25 111 V
-25 113 V
-25 114 V
-25 116 V
4195 2235 M
-25 80 V
-25 81 V
-25 81 V
-25 81 V
-25 82 V
-25 83 V
-25 84 V
-25 85 V
-26 86 V
-25 88 V
-25 89 V
-25 91 V
-25 92 V
-25 95 V
-25 97 V
-25 100 V
-25 101 V
-25 104 V
-25 107 V
-25 108 V
-25 110 V
-25 112 V
-25 114 V
-25 115 V
4313 2372 M
-25 80 V
-25 81 V
-25 82 V
-25 82 V
-25 83 V
-26 84 V
-25 84 V
-25 86 V
-25 87 V
-25 88 V
-25 90 V
-25 91 V
-25 94 V
-25 95 V
-25 97 V
-25 100 V
-25 101 V
-25 104 V
-25 106 V
-25 108 V
-25 109 V
-25 112 V
-25 113 V
-25 114 V
4431 2533 M
-25 81 V
-25 82 V
-25 82 V
-26 83 V
-25 83 V
-25 85 V
-25 85 V
-25 86 V
-25 88 V
-25 89 V
-25 90 V
-25 92 V
-25 94 V
-25 96 V
-25 97 V
-25 100 V
-25 101 V
-25 104 V
-25 105 V
-25 108 V
-25 109 V
-25 110 V
-26 113 V
-25 113 V
4549 2721 M
-26 81 V
-25 83 V
-25 82 V
-25 84 V
-25 84 V
-25 85 V
-25 86 V
-25 87 V
-25 88 V
-25 90 V
-25 91 V
-25 92 V
-25 95 V
-25 95 V
-25 98 V
-25 100 V
-25 101 V
-25 104 V
-25 105 V
-26 107 V
-25 108 V
-25 110 V
-25 112 V
stroke 3972 4889 M
-25 113 V
stroke
LT0
0.00 0.00 0.00 C 3960 4946 M
stroke
LT0
0.00 0.00 0.00 C 3960 4946 M
-21 42 V
47 -158 R
stroke
LT0
0.00 0.00 0.00 C 3986 4830 M
-22 45 V
stroke
LT0
0.00 0.00 0.00 C 3964 4875 M
-42 86 V
91 -253 R
stroke
LT0
0.00 0.00 0.00 C 4013 4708 M
-25 54 V
stroke
LT0
0.00 0.00 0.00 C 3988 4762 M
-41 85 V
stroke
LT0
0.00 0.00 0.00 C 3947 4847 M
-42 86 V
138 -355 R
stroke
LT0
0.00 0.00 0.00 C 4043 4578 M
-32 71 V
stroke
LT0
0.00 0.00 0.00 C 4011 4649 M
-40 85 V
stroke
LT0
0.00 0.00 0.00 C 3971 4734 M
-41 86 V
stroke
LT0
0.00 0.00 0.00 C 3930 4820 M
-42 85 V
188 -469 R
stroke
LT0
0.00 0.00 0.00 C 4076 4436 M
-5 13 V
stroke
LT0
0.00 0.00 0.00 C 4071 4449 M
-38 86 V
stroke
LT0
0.00 0.00 0.00 C 4033 4535 M
-39 86 V
stroke
LT0
0.00 0.00 0.00 C 3994 4621 M
-40 86 V
stroke
LT0
0.00 0.00 0.00 C 3954 4707 M
-41 85 V
stroke
LT0
0.00 0.00 0.00 C 3913 4792 M
-42 85 V
244 -601 R
stroke
LT0
0.00 0.00 0.00 C 4115 4276 M
-24 59 V
stroke
LT0
0.00 0.00 0.00 C 4091 4335 M
-37 86 V
stroke
LT0
0.00 0.00 0.00 C 4054 4421 M
-38 86 V
stroke
LT0
0.00 0.00 0.00 C 4016 4507 M
-38 86 V
stroke
LT0
0.00 0.00 0.00 C 3978 4593 M
-40 86 V
stroke
LT0
0.00 0.00 0.00 C 3938 4679 M
-42 85 V
stroke
LT0
0.00 0.00 0.00 C 3896 4764 M
-42 86 V
308 -766 R
stroke
LT0
0.00 0.00 0.00 C 4162 4084 M
-19 50 V
stroke
LT0
0.00 0.00 0.00 C 4143 4134 M
-34 87 V
stroke
LT0
0.00 0.00 0.00 C 4109 4221 M
-35 86 V
stroke
LT0
0.00 0.00 0.00 C 4074 4307 M
-36 87 V
stroke
LT0
0.00 0.00 0.00 C 4038 4394 M
-38 86 V
stroke
LT0
0.00 0.00 0.00 C 4000 4480 M
-39 86 V
stroke
LT0
0.00 0.00 0.00 C 3961 4566 M
-40 85 V
stroke
LT0
0.00 0.00 0.00 C 3921 4651 M
-41 86 V
stroke
LT0
0.00 0.00 0.00 C 3880 4737 M
-43 85 V
4233 3810 M
stroke
LT0
0.00 0.00 0.00 C 4233 3810 M
-12 36 V
stroke
LT0
0.00 0.00 0.00 C 4221 3846 M
-30 87 V
stroke
LT0
0.00 0.00 0.00 C 4191 3933 M
-32 87 V
stroke
LT0
0.00 0.00 0.00 C 4159 4020 M
-32 87 V
stroke
LT0
0.00 0.00 0.00 C 4127 4107 M
-34 86 V
stroke
LT0
0.00 0.00 0.00 C 4093 4193 M
-36 87 V
stroke
LT0
0.00 0.00 0.00 C 4057 4280 M
-36 86 V
stroke
LT0
0.00 0.00 0.00 C 4021 4366 M
-38 86 V
stroke
LT0
0.00 0.00 0.00 C 3983 4452 M
-39 86 V
stroke
LT0
0.00 0.00 0.00 C 3944 4538 M
-40 86 V
stroke
LT0
0.00 0.00 0.00 C 3904 4624 M
-41 85 V
stroke
LT0
0.00 0.00 0.00 C 3863 4709 M
-21 41 V
stroke
LT0
0.00 0.00 0.00 C 3842 4750 M
-23 44 V
4535 2700 M
stroke
LT0
0.00 0.00 0.00 C 4535 2700 M
-21 89 V
stroke
LT0
0.00 0.00 0.00 C 4514 2789 M
-21 89 V
stroke
LT0
0.00 0.00 0.00 C 4493 2878 M
-21 87 V
stroke
LT0
0.00 0.00 0.00 C 4472 2965 M
-9 36 V
54 -329 R
stroke
LT0
0.00 0.00 0.00 C 4517 2672 M
-20 89 V
stroke
LT0
0.00 0.00 0.00 C 4497 2761 M
-21 89 V
stroke
LT0
0.00 0.00 0.00 C 4476 2850 M
-21 88 V
stroke
LT0
0.00 0.00 0.00 C 4455 2938 M
-22 88 V
stroke
LT0
0.00 0.00 0.00 C 4433 3026 M
-22 89 V
stroke
LT0
0.00 0.00 0.00 C 4411 3115 M
-23 88 V
stroke
LT0
0.00 0.00 0.00 C 4388 3203 M
-24 88 V
stroke
LT0
0.00 0.00 0.00 C 4364 3291 M
-24 88 V
stroke
LT0
0.00 0.00 0.00 C 4340 3379 M
-25 88 V
stroke
LT0
0.00 0.00 0.00 C 4315 3467 M
-26 88 V
stroke
LT0
0.00 0.00 0.00 C 4289 3555 M
-27 88 V
stroke
LT0
0.00 0.00 0.00 C 4262 3643 M
-28 88 V
stroke
LT0
0.00 0.00 0.00 C 4234 3731 M
-29 87 V
stroke
LT0
0.00 0.00 0.00 C 4205 3818 M
-31 87 V
stroke
LT0
0.00 0.00 0.00 C 4174 3905 M
-31 87 V
stroke
LT0
0.00 0.00 0.00 C 4143 3992 M
-33 87 V
stroke
LT0
0.00 0.00 0.00 C 4110 4079 M
-34 87 V
stroke
LT0
0.00 0.00 0.00 C 4076 4166 M
-35 86 V
stroke
LT0
0.00 0.00 0.00 C 4041 4252 M
-37 86 V
stroke
LT0
0.00 0.00 0.00 C 4004 4338 M
-38 86 V
stroke
LT0
0.00 0.00 0.00 C 3966 4424 M
-39 86 V
stroke
LT0
0.00 0.00 0.00 C 3927 4510 M
-40 86 V
stroke
LT0
0.00 0.00 0.00 C 3887 4596 M
-19 40 V
stroke
LT0
0.00 0.00 0.00 C 3868 4636 M
-24 45 V
stroke
LT0
0.00 0.00 0.00 C 3844 4681 M
-45 85 V
4500 2645 M
stroke
LT0
0.00 0.00 0.00 C 4500 2645 M
-20 88 V
stroke
LT0
0.00 0.00 0.00 C 4480 2733 M
-21 89 V
stroke
LT0
0.00 0.00 0.00 C 4459 2822 M
-21 89 V
stroke
LT0
0.00 0.00 0.00 C 4438 2911 M
-22 87 V
stroke
LT0
0.00 0.00 0.00 C 4416 2998 M
-22 89 V
stroke
LT0
0.00 0.00 0.00 C 4394 3087 M
-23 88 V
stroke
LT0
0.00 0.00 0.00 C 4371 3175 M
-24 89 V
stroke
LT0
0.00 0.00 0.00 C 4347 3264 M
-24 88 V
stroke
LT0
0.00 0.00 0.00 C 4323 3352 M
-25 88 V
stroke
LT0
0.00 0.00 0.00 C 4298 3440 M
-26 88 V
stroke
LT0
0.00 0.00 0.00 C 4272 3528 M
-27 87 V
stroke
LT0
0.00 0.00 0.00 C 4245 3615 M
-28 88 V
stroke
LT0
0.00 0.00 0.00 C 4217 3703 M
-29 87 V
stroke
LT0
0.00 0.00 0.00 C 4188 3790 M
-30 88 V
stroke
LT0
0.00 0.00 0.00 C 4158 3878 M
-32 87 V
stroke
LT0
0.00 0.00 0.00 C 4126 3965 M
-33 87 V
stroke
LT0
0.00 0.00 0.00 C 4093 4052 M
-34 86 V
stroke
LT0
0.00 0.00 0.00 C 4059 4138 M
-35 87 V
stroke
LT0
0.00 0.00 0.00 C 4024 4225 M
-36 86 V
stroke
LT0
0.00 0.00 0.00 C 3988 4311 M
-38 86 V
stroke
LT0
0.00 0.00 0.00 C 3950 4397 M
-39 86 V
stroke
LT0
0.00 0.00 0.00 C 3911 4483 M
-16 34 V
stroke
LT0
0.00 0.00 0.00 C 3895 4517 M
-26 51 V
stroke
LT0
0.00 0.00 0.00 C 3869 4568 M
-45 85 V
stroke
LT0
0.00 0.00 0.00 C 3824 4653 M
-45 85 V
4483 2617 M
stroke
LT0
0.00 0.00 0.00 C 4483 2617 M
-21 89 V
stroke
LT0
0.00 0.00 0.00 C 4462 2706 M
-20 88 V
stroke
LT0
0.00 0.00 0.00 C 4442 2794 M
-22 89 V
stroke
LT0
0.00 0.00 0.00 C 4420 2883 M
-21 88 V
stroke
LT0
0.00 0.00 0.00 C 4399 2971 M
-23 88 V
stroke
LT0
0.00 0.00 0.00 C 4376 3059 M
-22 89 V
stroke
LT0
0.00 0.00 0.00 C 4354 3148 M
-24 88 V
stroke
LT0
0.00 0.00 0.00 C 4330 3236 M
-24 88 V
stroke
LT0
0.00 0.00 0.00 C 4306 3324 M
-25 88 V
stroke
LT0
0.00 0.00 0.00 C 4281 3412 M
-26 88 V
stroke
LT0
0.00 0.00 0.00 C 4255 3500 M
-27 88 V
stroke
LT0
0.00 0.00 0.00 C 4228 3588 M
-28 87 V
stroke
LT0
0.00 0.00 0.00 C 4200 3675 M
-29 88 V
stroke
LT0
0.00 0.00 0.00 C 4171 3763 M
-30 87 V
stroke
LT0
0.00 0.00 0.00 C 4141 3850 M
-31 87 V
stroke
LT0
0.00 0.00 0.00 C 4110 3937 M
-33 87 V
stroke
LT0
0.00 0.00 0.00 C 4077 4024 M
-34 87 V
stroke
LT0
0.00 0.00 0.00 C 4043 4111 M
-35 86 V
stroke
LT0
0.00 0.00 0.00 C 4008 4197 M
-37 86 V
stroke
LT0
0.00 0.00 0.00 C 3971 4283 M
-38 86 V
stroke
LT0
0.00 0.00 0.00 C 3933 4369 M
-9 21 V
stroke
LT0
0.00 0.00 0.00 C 3924 4390 M
-32 65 V
stroke
LT0
0.00 0.00 0.00 C 3892 4455 M
-43 85 V
stroke
LT0
0.00 0.00 0.00 C 3849 4540 M
-44 85 V
stroke
LT0
0.00 0.00 0.00 C 3805 4625 M
-46 85 V
4465 2589 M
stroke
LT0
0.00 0.00 0.00 C 4465 2589 M
-20 89 V
stroke
LT0
0.00 0.00 0.00 C 4445 2678 M
-21 89 V
stroke
LT0
0.00 0.00 0.00 C 4424 2767 M
-21 88 V
stroke
LT0
0.00 0.00 0.00 C 4403 2855 M
-22 88 V
stroke
LT0
0.00 0.00 0.00 C 4381 2943 M
-22 88 V
stroke
LT0
0.00 0.00 0.00 C 4359 3031 M
-22 89 V
stroke
LT0
0.00 0.00 0.00 C 4337 3120 M
-24 88 V
stroke
LT0
0.00 0.00 0.00 C 4313 3208 M
-24 88 V
stroke
LT0
0.00 0.00 0.00 C 4289 3296 M
-25 88 V
stroke
LT0
0.00 0.00 0.00 C 4264 3384 M
-26 88 V
stroke
LT0
0.00 0.00 0.00 C 4238 3472 M
-27 88 V
stroke
LT0
0.00 0.00 0.00 C 4211 3560 M
-27 88 V
stroke
LT0
0.00 0.00 0.00 C 4184 3648 M
-29 87 V
stroke
LT0
0.00 0.00 0.00 C 4155 3735 M
-31 87 V
stroke
LT0
0.00 0.00 0.00 C 4124 3822 M
-31 87 V
stroke
LT0
0.00 0.00 0.00 C 4093 3909 M
-33 87 V
stroke
LT0
0.00 0.00 0.00 C 4060 3996 M
-34 87 V
stroke
LT0
0.00 0.00 0.00 C 4026 4083 M
-35 86 V
stroke
LT0
0.00 0.00 0.00 C 3991 4169 M
-35 83 V
stroke
LT0
0.00 0.00 0.00 C 3956 4252 M
-2 4 V
stroke
LT0
0.00 0.00 0.00 C 3954 4256 M
-40 85 V
stroke
LT0
0.00 0.00 0.00 C 3914 4341 M
-42 86 V
stroke
LT0
0.00 0.00 0.00 C 3872 4427 M
-43 85 V
stroke
LT0
0.00 0.00 0.00 C 3829 4512 M
-44 85 V
stroke
LT0
0.00 0.00 0.00 C 3785 4597 M
-46 85 V
4448 2561 M
stroke
LT0
0.00 0.00 0.00 C 4448 2561 M
-20 89 V
stroke
LT0
0.00 0.00 0.00 C 4428 2650 M
-21 89 V
stroke
LT0
0.00 0.00 0.00 C 4407 2739 M
-21 89 V
stroke
LT0
0.00 0.00 0.00 C 4386 2828 M
-22 88 V
stroke
LT0
0.00 0.00 0.00 C 4364 2916 M
-22 88 V
stroke
LT0
0.00 0.00 0.00 C 4342 3004 M
-22 88 V
stroke
LT0
0.00 0.00 0.00 C 4320 3092 M
-24 88 V
stroke
LT0
0.00 0.00 0.00 C 4296 3180 M
-24 89 V
stroke
LT0
0.00 0.00 0.00 C 4272 3269 M
-25 88 V
stroke
LT0
0.00 0.00 0.00 C 4247 3357 M
-26 88 V
stroke
LT0
0.00 0.00 0.00 C 4221 3445 M
-26 87 V
stroke
LT0
0.00 0.00 0.00 C 4195 3532 M
-28 88 V
stroke
LT0
0.00 0.00 0.00 C 4167 3620 M
-29 88 V
stroke
LT0
0.00 0.00 0.00 C 4138 3708 M
-30 87 V
stroke
LT0
0.00 0.00 0.00 C 4108 3795 M
-32 87 V
stroke
LT0
0.00 0.00 0.00 C 4076 3882 M
-32 87 V
stroke
LT0
0.00 0.00 0.00 C 4044 3969 M
-34 86 V
stroke
LT0
0.00 0.00 0.00 C 4010 4055 M
-17 41 V
stroke
LT0
0.00 0.00 0.00 C 3993 4096 M
-20 46 V
stroke
LT0
0.00 0.00 0.00 C 3973 4142 M
-38 85 V
stroke
LT0
0.00 0.00 0.00 C 3935 4227 M
-41 86 V
stroke
LT0
0.00 0.00 0.00 C 3894 4313 M
-41 85 V
stroke
LT0
0.00 0.00 0.00 C 3853 4398 M
-44 86 V
stroke
LT0
0.00 0.00 0.00 C 3809 4484 M
-44 85 V
stroke
LT0
0.00 0.00 0.00 C 3765 4569 M
-46 84 V
4430 2534 M
stroke
LT0
0.00 0.00 0.00 C 4430 2534 M
-1 6 V
stroke
LT0
0.00 0.00 0.00 C 4429 2540 M
-19 82 V
stroke
LT0
0.00 0.00 0.00 C 4410 2622 M
-20 89 V
stroke
LT0
0.00 0.00 0.00 C 4390 2711 M
-21 89 V
stroke
LT0
0.00 0.00 0.00 C 4369 2800 M
-22 88 V
stroke
LT0
0.00 0.00 0.00 C 4347 2888 M
-22 88 V
stroke
LT0
0.00 0.00 0.00 C 4325 2976 M
-23 88 V
stroke
LT0
0.00 0.00 0.00 C 4302 3064 M
-23 89 V
stroke
LT0
0.00 0.00 0.00 C 4279 3153 M
-24 88 V
stroke
LT0
0.00 0.00 0.00 C 4255 3241 M
-25 88 V
stroke
LT0
0.00 0.00 0.00 C 4230 3329 M
-25 88 V
stroke
LT0
0.00 0.00 0.00 C 4205 3417 M
-27 88 V
stroke
LT0
0.00 0.00 0.00 C 4178 3505 M
-28 87 V
stroke
LT0
0.00 0.00 0.00 C 4150 3592 M
-29 88 V
stroke
LT0
0.00 0.00 0.00 C 4121 3680 M
-30 87 V
stroke
LT0
0.00 0.00 0.00 C 4091 3767 M
-31 87 V
stroke
LT0
0.00 0.00 0.00 C 4060 3854 M
-21 56 V
stroke
LT0
0.00 0.00 0.00 C 4039 3910 M
-12 31 V
stroke
LT0
0.00 0.00 0.00 C 4027 3941 M
-36 86 V
stroke
LT0
0.00 0.00 0.00 C 3991 4027 M
-37 86 V
stroke
LT0
0.00 0.00 0.00 C 3954 4113 M
-39 86 V
stroke
LT0
0.00 0.00 0.00 C 3915 4199 M
-40 86 V
stroke
LT0
0.00 0.00 0.00 C 3875 4285 M
-42 85 V
stroke
LT0
0.00 0.00 0.00 C 3833 4370 M
-43 85 V
stroke
LT0
0.00 0.00 0.00 C 3790 4455 M
-45 85 V
stroke
LT0
0.00 0.00 0.00 C 3745 4540 M
-19 35 V
stroke
LT0
0.00 0.00 0.00 C 3726 4575 M
-29 50 V
4410 2505 M
stroke
LT0
0.00 0.00 0.00 C 4410 2505 M
-20 89 V
stroke
LT0
0.00 0.00 0.00 C 4390 2594 M
-19 89 V
stroke
LT0
0.00 0.00 0.00 C 4371 2683 M
-20 89 V
stroke
LT0
0.00 0.00 0.00 C 4351 2772 M
-21 89 V
stroke
LT0
0.00 0.00 0.00 C 4330 2861 M
-4 15 V
stroke
LT0
0.00 0.00 0.00 C 4326 2876 M
-18 72 V
stroke
LT0
0.00 0.00 0.00 C 4308 2948 M
-23 89 V
stroke
LT0
0.00 0.00 0.00 C 4285 3037 M
-23 88 V
stroke
LT0
0.00 0.00 0.00 C 4262 3125 M
-24 88 V
stroke
LT0
0.00 0.00 0.00 C 4238 3213 M
-25 88 V
stroke
LT0
0.00 0.00 0.00 C 4213 3301 M
-25 88 V
stroke
LT0
0.00 0.00 0.00 C 4188 3389 M
-27 88 V
stroke
LT0
0.00 0.00 0.00 C 4161 3477 M
-28 88 V
stroke
LT0
0.00 0.00 0.00 C 4133 3565 M
-25 78 V
stroke
LT0
0.00 0.00 0.00 C 4108 3643 M
-4 9 V
stroke
LT0
0.00 0.00 0.00 C 4104 3652 M
-30 87 V
stroke
LT0
0.00 0.00 0.00 C 4074 3739 M
-33 87 V
stroke
LT0
0.00 0.00 0.00 C 4041 3826 M
-34 87 V
stroke
LT0
0.00 0.00 0.00 C 4007 3913 M
-35 86 V
stroke
LT0
0.00 0.00 0.00 C 3972 3999 M
-38 86 V
stroke
LT0
0.00 0.00 0.00 C 3934 4085 M
-39 86 V
stroke
LT0
0.00 0.00 0.00 C 3895 4171 M
-40 86 V
stroke
LT0
0.00 0.00 0.00 C 3855 4257 M
-42 85 V
stroke
LT0
0.00 0.00 0.00 C 3813 4342 M
-43 85 V
stroke
LT0
0.00 0.00 0.00 C 3770 4427 M
-19 36 V
stroke
LT0
0.00 0.00 0.00 C 3751 4463 M
-28 49 V
stroke
LT0
0.00 0.00 0.00 C 3723 4512 M
-50 84 V
4389 2477 M
stroke
LT0
0.00 0.00 0.00 C 4389 2477 M
-19 89 V
stroke
LT0
0.00 0.00 0.00 C 4370 2566 M
-20 89 V
stroke
LT0
0.00 0.00 0.00 C 4350 2655 M
-20 89 V
stroke
LT0
0.00 0.00 0.00 C 4330 2744 M
-21 89 V
stroke
LT0
0.00 0.00 0.00 C 4309 2833 M
-21 88 V
stroke
LT0
0.00 0.00 0.00 C 4288 2921 M
-22 88 V
stroke
LT0
0.00 0.00 0.00 C 4266 3009 M
-23 88 V
stroke
LT0
0.00 0.00 0.00 C 4243 3097 M
-23 88 V
stroke
LT0
0.00 0.00 0.00 C 4220 3185 M
-25 89 V
stroke
LT0
0.00 0.00 0.00 C 4195 3274 M
-26 87 V
stroke
LT0
0.00 0.00 0.00 C 4169 3361 M
-26 88 V
stroke
LT0
0.00 0.00 0.00 C 4143 3449 M
-29 88 V
stroke
LT0
0.00 0.00 0.00 C 4114 3537 M
-29 87 V
stroke
LT0
0.00 0.00 0.00 C 4085 3624 M
-31 87 V
stroke
LT0
0.00 0.00 0.00 C 4054 3711 M
-32 87 V
stroke
LT0
0.00 0.00 0.00 C 4022 3798 M
-34 87 V
stroke
LT0
0.00 0.00 0.00 C 3988 3885 M
-36 86 V
stroke
LT0
0.00 0.00 0.00 C 3952 3971 M
-37 86 V
stroke
LT0
0.00 0.00 0.00 C 3915 4057 M
-39 86 V
stroke
LT0
0.00 0.00 0.00 C 3876 4143 M
-41 86 V
stroke
LT0
0.00 0.00 0.00 C 3835 4229 M
-42 85 V
stroke
LT0
0.00 0.00 0.00 C 3793 4314 M
-16 32 V
stroke
LT0
0.00 0.00 0.00 C 3777 4346 M
-29 53 V
stroke
LT0
0.00 0.00 0.00 C 3748 4399 M
-49 84 V
stroke
LT0
0.00 0.00 0.00 C 3699 4483 M
-50 84 V
4368 2449 M
stroke
LT0
0.00 0.00 0.00 C 4368 2449 M
-19 89 V
stroke
LT0
0.00 0.00 0.00 C 4349 2538 M
-19 89 V
stroke
LT0
0.00 0.00 0.00 C 4330 2627 M
-20 89 V
stroke
LT0
0.00 0.00 0.00 C 4310 2716 M
-21 88 V
stroke
LT0
0.00 0.00 0.00 C 4289 2804 M
-21 89 V
stroke
LT0
0.00 0.00 0.00 C 4268 2893 M
-22 88 V
stroke
LT0
0.00 0.00 0.00 C 4246 2981 M
-23 88 V
stroke
LT0
0.00 0.00 0.00 C 4223 3069 M
-23 88 V
stroke
LT0
0.00 0.00 0.00 C 4200 3157 M
-25 88 V
stroke
LT0
0.00 0.00 0.00 C 4175 3245 M
-25 88 V
stroke
LT0
0.00 0.00 0.00 C 4150 3333 M
-27 88 V
stroke
LT0
0.00 0.00 0.00 C 4123 3421 M
-28 88 V
stroke
LT0
0.00 0.00 0.00 C 4095 3509 M
-30 87 V
stroke
LT0
0.00 0.00 0.00 C 4065 3596 M
-30 87 V
stroke
LT0
0.00 0.00 0.00 C 4035 3683 M
-33 87 V
stroke
LT0
0.00 0.00 0.00 C 4002 3770 M
-34 87 V
stroke
LT0
0.00 0.00 0.00 C 3968 3857 M
-35 86 V
stroke
LT0
0.00 0.00 0.00 C 3933 3943 M
-38 86 V
stroke
LT0
0.00 0.00 0.00 C 3895 4029 M
-39 86 V
stroke
LT0
0.00 0.00 0.00 C 3856 4115 M
-40 86 V
stroke
LT0
0.00 0.00 0.00 C 3816 4201 M
-10 20 V
stroke
LT0
0.00 0.00 0.00 C 3806 4221 M
-35 65 V
stroke
LT0
0.00 0.00 0.00 C 3771 4286 M
-47 84 V
stroke
LT0
0.00 0.00 0.00 C 3724 4370 M
-49 84 V
stroke
LT0
0.00 0.00 0.00 C 3675 4454 M
-50 84 V
4348 2420 M
stroke
LT0
0.00 0.00 0.00 C 4348 2420 M
-19 89 V
stroke
LT0
0.00 0.00 0.00 C 4329 2509 M
-20 89 V
stroke
LT0
0.00 0.00 0.00 C 4309 2598 M
-20 89 V
stroke
LT0
0.00 0.00 0.00 C 4289 2687 M
-20 89 V
stroke
LT0
0.00 0.00 0.00 C 4269 2776 M
-21 89 V
stroke
LT0
0.00 0.00 0.00 C 4248 2865 M
-22 87 V
stroke
LT0
0.00 0.00 0.00 C 4226 2952 M
-23 89 V
stroke
LT0
0.00 0.00 0.00 C 4203 3041 M
-23 88 V
stroke
LT0
0.00 0.00 0.00 C 4180 3129 M
-25 88 V
stroke
LT0
0.00 0.00 0.00 C 4155 3217 M
-25 88 V
stroke
LT0
0.00 0.00 0.00 C 4130 3305 M
-27 88 V
stroke
LT0
0.00 0.00 0.00 C 4103 3393 M
-28 88 V
stroke
LT0
0.00 0.00 0.00 C 4075 3481 M
-29 87 V
stroke
LT0
0.00 0.00 0.00 C 4046 3568 M
-31 87 V
stroke
LT0
0.00 0.00 0.00 C 4015 3655 M
-32 87 V
stroke
LT0
0.00 0.00 0.00 C 3983 3742 M
-34 87 V
stroke
LT0
0.00 0.00 0.00 C 3949 3829 M
-36 86 V
stroke
LT0
0.00 0.00 0.00 C 3913 3915 M
-37 86 V
stroke
LT0
0.00 0.00 0.00 C 3876 4001 M
-39 86 V
stroke
LT0
0.00 0.00 0.00 C 3837 4087 M
stroke
LT0
0.00 0.00 0.00 C 3837 4087 M
-44 85 V
stroke
LT0
0.00 0.00 0.00 C 3793 4172 M
-45 85 V
stroke
LT0
0.00 0.00 0.00 C 3748 4257 M
-47 84 V
stroke
LT0
0.00 0.00 0.00 C 3701 4341 M
-49 85 V
stroke
LT0
0.00 0.00 0.00 C 3652 4426 M
-51 84 V
4327 2392 M
stroke
LT0
0.00 0.00 0.00 C 4327 2392 M
-19 89 V
stroke
LT0
0.00 0.00 0.00 C 4308 2481 M
-19 89 V
stroke
LT0
0.00 0.00 0.00 C 4289 2570 M
-20 89 V
stroke
LT0
0.00 0.00 0.00 C 4269 2659 M
-21 89 V
stroke
LT0
0.00 0.00 0.00 C 4248 2748 M
-21 89 V
stroke
LT0
0.00 0.00 0.00 C 4227 2837 M
-21 88 V
stroke
LT0
0.00 0.00 0.00 C 4206 2925 M
-23 88 V
stroke
LT0
0.00 0.00 0.00 C 4183 3013 M
-23 88 V
stroke
LT0
0.00 0.00 0.00 C 4160 3101 M
-25 88 V
stroke
LT0
0.00 0.00 0.00 C 4135 3189 M
-25 88 V
stroke
LT0
0.00 0.00 0.00 C 4110 3277 M
-27 88 V
stroke
LT0
0.00 0.00 0.00 C 4083 3365 M
-27 87 V
stroke
LT0
0.00 0.00 0.00 C 4056 3452 M
-30 88 V
stroke
LT0
0.00 0.00 0.00 C 4026 3540 M
-31 87 V
stroke
LT0
0.00 0.00 0.00 C 3995 3627 M
-32 87 V
stroke
LT0
0.00 0.00 0.00 C 3963 3714 M
-34 87 V
stroke
LT0
0.00 0.00 0.00 C 3929 3801 M
-35 86 V
stroke
LT0
0.00 0.00 0.00 C 3894 3887 M
-21 47 V
stroke
LT0
0.00 0.00 0.00 C 3873 3934 M
-18 39 V
stroke
LT0
0.00 0.00 0.00 C 3855 3973 M
-41 85 V
stroke
LT0
0.00 0.00 0.00 C 3814 4058 M
-44 85 V
stroke
LT0
0.00 0.00 0.00 C 3770 4143 M
-46 85 V
stroke
LT0
0.00 0.00 0.00 C 3724 4228 M
-47 85 V
stroke
LT0
0.00 0.00 0.00 C 3677 4313 M
-49 84 V
stroke
LT0
0.00 0.00 0.00 C 3628 4397 M
-18 30 V
stroke
LT0
0.00 0.00 0.00 C 3610 4427 M
-37 53 V
4305 2364 M
stroke
LT0
0.00 0.00 0.00 C 4305 2364 M
-18 89 V
stroke
LT0
0.00 0.00 0.00 C 4287 2453 M
-1 6 V
stroke
LT0
0.00 0.00 0.00 C 4286 2459 M
-18 83 V
stroke
LT0
0.00 0.00 0.00 C 4268 2542 M
-20 89 V
stroke
LT0
0.00 0.00 0.00 C 4248 2631 M
-20 89 V
stroke
LT0
0.00 0.00 0.00 C 4228 2720 M
-21 88 V
stroke
LT0
0.00 0.00 0.00 C 4207 2808 M
-21 89 V
stroke
LT0
0.00 0.00 0.00 C 4186 2897 M
-23 87 V
stroke
LT0
0.00 0.00 0.00 C 4163 2984 M
-23 89 V
stroke
LT0
0.00 0.00 0.00 C 4140 3073 M
-24 88 V
stroke
LT0
0.00 0.00 0.00 C 4116 3161 M
-26 88 V
stroke
LT0
0.00 0.00 0.00 C 4090 3249 M
-26 88 V
stroke
LT0
0.00 0.00 0.00 C 4064 3337 M
-28 87 V
stroke
LT0
0.00 0.00 0.00 C 4036 3424 M
-29 88 V
stroke
LT0
0.00 0.00 0.00 C 4007 3512 M
-31 87 V
stroke
LT0
0.00 0.00 0.00 C 3976 3599 M
-32 87 V
stroke
LT0
0.00 0.00 0.00 C 3944 3686 M
-26 65 V
stroke
LT0
0.00 0.00 0.00 C 3918 3751 M
-9 21 V
stroke
LT0
0.00 0.00 0.00 C 3909 3772 M
-37 86 V
stroke
LT0
0.00 0.00 0.00 C 3872 3858 M
-40 86 V
stroke
LT0
0.00 0.00 0.00 C 3832 3944 M
-42 86 V
stroke
LT0
0.00 0.00 0.00 C 3790 4030 M
-43 85 V
stroke
LT0
0.00 0.00 0.00 C 3747 4115 M
-46 84 V
stroke
LT0
0.00 0.00 0.00 C 3701 4199 M
-48 85 V
stroke
LT0
0.00 0.00 0.00 C 3653 4284 M
-19 33 V
stroke
LT0
0.00 0.00 0.00 C 3634 4317 M
-33 51 V
stroke
LT0
0.00 0.00 0.00 C 3601 4368 M
-57 83 V
4280 2335 M
stroke
LT0
0.00 0.00 0.00 C 4280 2335 M
-18 89 V
stroke
LT0
0.00 0.00 0.00 C 4262 2424 M
-17 89 V
stroke
LT0
0.00 0.00 0.00 C 4245 2513 M
-19 89 V
stroke
LT0
0.00 0.00 0.00 C 4226 2602 M
-19 89 V
stroke
LT0
0.00 0.00 0.00 C 4207 2691 M
-20 89 V
stroke
LT0
0.00 0.00 0.00 C 4187 2780 M
-4 16 V
stroke
LT0
0.00 0.00 0.00 C 4183 2796 M
-18 73 V
stroke
LT0
0.00 0.00 0.00 C 4165 2869 M
-22 87 V
stroke
LT0
0.00 0.00 0.00 C 4143 2956 M
-23 89 V
stroke
LT0
0.00 0.00 0.00 C 4120 3045 M
-24 88 V
stroke
LT0
0.00 0.00 0.00 C 4096 3133 M
-26 88 V
stroke
LT0
0.00 0.00 0.00 C 4070 3221 M
-26 88 V
stroke
LT0
0.00 0.00 0.00 C 4044 3309 M
-28 87 V
stroke
LT0
0.00 0.00 0.00 C 4016 3396 M
-29 88 V
stroke
LT0
0.00 0.00 0.00 C 3987 3484 M
stroke
LT0
0.00 0.00 0.00 C 3987 3484 M
-32 87 V
stroke
LT0
0.00 0.00 0.00 C 3955 3571 M
-33 86 V
stroke
LT0
0.00 0.00 0.00 C 3922 3657 M
-36 87 V
stroke
LT0
0.00 0.00 0.00 C 3886 3744 M
-37 86 V
stroke
LT0
0.00 0.00 0.00 C 3849 3830 M
-40 85 V
stroke
LT0
0.00 0.00 0.00 C 3809 3915 M
-42 86 V
stroke
LT0
0.00 0.00 0.00 C 3767 4001 M
-44 85 V
stroke
LT0
0.00 0.00 0.00 C 3723 4086 M
-46 85 V
stroke
LT0
0.00 0.00 0.00 C 3677 4171 M
-18 32 V
stroke
LT0
0.00 0.00 0.00 C 3659 4203 M
-32 52 V
stroke
LT0
0.00 0.00 0.00 C 3627 4255 M
-55 83 V
stroke
LT0
0.00 0.00 0.00 C 3572 4338 M
-57 83 V
4255 2306 M
stroke
LT0
0.00 0.00 0.00 C 4255 2306 M
-17 89 V
stroke
LT0
0.00 0.00 0.00 C 4238 2395 M
-18 89 V
stroke
LT0
0.00 0.00 0.00 C 4220 2484 M
-19 89 V
stroke
LT0
0.00 0.00 0.00 C 4201 2573 M
-19 89 V
stroke
LT0
0.00 0.00 0.00 C 4182 2662 M
-20 89 V
stroke
LT0
0.00 0.00 0.00 C 4162 2751 M
-20 89 V
stroke
LT0
0.00 0.00 0.00 C 4142 2840 M
-22 89 V
stroke
LT0
0.00 0.00 0.00 C 4120 2929 M
-23 87 V
stroke
LT0
0.00 0.00 0.00 C 4097 3016 M
-23 88 V
stroke
LT0
0.00 0.00 0.00 C 4074 3104 M
-26 88 V
stroke
LT0
0.00 0.00 0.00 C 4048 3192 M
-26 88 V
stroke
LT0
0.00 0.00 0.00 C 4022 3280 M
-28 88 V
stroke
LT0
0.00 0.00 0.00 C 3994 3368 M
-30 87 V
stroke
LT0
0.00 0.00 0.00 C 3964 3455 M
-32 87 V
stroke
LT0
0.00 0.00 0.00 C 3932 3542 M
-33 87 V
stroke
LT0
0.00 0.00 0.00 C 3899 3629 M
-36 86 V
stroke
LT0
0.00 0.00 0.00 C 3863 3715 M
-38 86 V
stroke
LT0
0.00 0.00 0.00 C 3825 3801 M
-39 86 V
stroke
LT0
0.00 0.00 0.00 C 3786 3887 M
-42 85 V
stroke
LT0
0.00 0.00 0.00 C 3744 3972 M
-44 85 V
stroke
LT0
0.00 0.00 0.00 C 3700 4057 M
-13 24 V
stroke
LT0
0.00 0.00 0.00 C 3687 4081 M
-36 60 V
stroke
LT0
0.00 0.00 0.00 C 3651 4141 M
-53 84 V
stroke
LT0
0.00 0.00 0.00 C 3598 4225 M
-55 83 V
stroke
LT0
0.00 0.00 0.00 C 3543 4308 M
-57 83 V
4230 2277 M
stroke
LT0
0.00 0.00 0.00 C 4230 2277 M
-17 89 V
stroke
LT0
0.00 0.00 0.00 C 4213 2366 M
-18 89 V
stroke
LT0
0.00 0.00 0.00 C 4195 2455 M
-18 89 V
stroke
LT0
0.00 0.00 0.00 C 4177 2544 M
-19 89 V
stroke
LT0
0.00 0.00 0.00 C 4158 2633 M
-20 89 V
stroke
LT0
0.00 0.00 0.00 C 4138 2722 M
-20 89 V
stroke
LT0
0.00 0.00 0.00 C 4118 2811 M
-22 89 V
stroke
LT0
0.00 0.00 0.00 C 4096 2900 M
-22 87 V
stroke
LT0
0.00 0.00 0.00 C 4074 2987 M
-24 88 V
stroke
LT0
0.00 0.00 0.00 C 4050 3075 M
-25 89 V
stroke
LT0
0.00 0.00 0.00 C 4025 3164 M
-27 87 V
stroke
LT0
0.00 0.00 0.00 C 3998 3251 M
-28 88 V
stroke
LT0
0.00 0.00 0.00 C 3970 3339 M
-29 87 V
stroke
LT0
0.00 0.00 0.00 C 3941 3426 M
-32 87 V
stroke
LT0
0.00 0.00 0.00 C 3909 3513 M
-34 87 V
stroke
LT0
0.00 0.00 0.00 C 3875 3600 M
-35 86 V
stroke
LT0
0.00 0.00 0.00 C 3840 3686 M
-38 86 V
stroke
LT0
0.00 0.00 0.00 C 3802 3772 M
-40 86 V
stroke
LT0
0.00 0.00 0.00 C 3762 3858 M
-42 85 V
stroke
LT0
0.00 0.00 0.00 C 3720 3943 M
-3 7 V
stroke
LT0
0.00 0.00 0.00 C 3717 3950 M
-44 78 V
stroke
LT0
0.00 0.00 0.00 C 3673 4028 M
-51 84 V
stroke
LT0
0.00 0.00 0.00 C 3622 4112 M
-53 84 V
stroke
LT0
0.00 0.00 0.00 C 3569 4196 M
-55 83 V
stroke
LT0
0.00 0.00 0.00 C 3514 4279 M
-24 34 V
stroke
LT0
0.00 0.00 0.00 C 3490 4313 M
-38 48 V
4205 2248 M
stroke
LT0
0.00 0.00 0.00 C 4205 2248 M
-17 89 V
stroke
LT0
0.00 0.00 0.00 C 4188 2337 M
-18 89 V
stroke
LT0
0.00 0.00 0.00 C 4170 2426 M
-18 90 V
stroke
LT0
0.00 0.00 0.00 C 4152 2516 M
-19 89 V
stroke
LT0
0.00 0.00 0.00 C 4133 2605 M
-19 89 V
stroke
LT0
0.00 0.00 0.00 C 4114 2694 M
-21 88 V
stroke
LT0
0.00 0.00 0.00 C 4093 2782 M
-21 89 V
stroke
LT0
0.00 0.00 0.00 C 4072 2871 M
-22 87 V
stroke
LT0
0.00 0.00 0.00 C 4050 2958 M
-24 89 V
stroke
LT0
0.00 0.00 0.00 C 4026 3047 M
-25 88 V
stroke
LT0
0.00 0.00 0.00 C 4001 3135 M
-26 88 V
stroke
LT0
0.00 0.00 0.00 C 3975 3223 M
-28 87 V
stroke
LT0
0.00 0.00 0.00 C 3947 3310 M
-30 88 V
stroke
LT0
0.00 0.00 0.00 C 3917 3398 M
-31 87 V
stroke
LT0
0.00 0.00 0.00 C 3886 3485 M
-34 86 V
stroke
LT0
0.00 0.00 0.00 C 3852 3571 M
-35 87 V
stroke
LT0
0.00 0.00 0.00 C 3817 3658 M
-38 86 V
stroke
LT0
0.00 0.00 0.00 C 3779 3744 M
-27 58 V
stroke
LT0
0.00 0.00 0.00 C 3752 3802 M
-14 27 V
stroke
LT0
0.00 0.00 0.00 C 3738 3829 M
-45 85 V
stroke
LT0
0.00 0.00 0.00 C 3693 3914 M
-49 84 V
stroke
LT0
0.00 0.00 0.00 C 3644 3998 M
-50 84 V
stroke
LT0
0.00 0.00 0.00 C 3594 4082 M
-53 84 V
stroke
LT0
0.00 0.00 0.00 C 3541 4166 M
-27 40 V
stroke
LT0
0.00 0.00 0.00 C 3514 4206 M
-33 43 V
stroke
LT0
0.00 0.00 0.00 C 3481 4249 M
-65 81 V
4177 2218 M
stroke
LT0
0.00 0.00 0.00 C 4177 2218 M
-16 90 V
stroke
LT0
0.00 0.00 0.00 C 4161 2308 M
-14 80 V
stroke
LT0
0.00 0.00 0.00 C 4147 2388 M
-1 9 V
stroke
LT0
0.00 0.00 0.00 C 4146 2397 M
-18 90 V
stroke
LT0
0.00 0.00 0.00 C 4128 2487 M
-19 89 V
stroke
LT0
0.00 0.00 0.00 C 4109 2576 M
-20 89 V
stroke
LT0
0.00 0.00 0.00 C 4089 2665 M
-20 88 V
stroke
LT0
0.00 0.00 0.00 C 4069 2753 M
-21 89 V
stroke
LT0
0.00 0.00 0.00 C 4048 2842 M
-22 89 V
stroke
LT0
0.00 0.00 0.00 C 4026 2931 M
-24 87 V
stroke
LT0
0.00 0.00 0.00 C 4002 3018 M
-24 88 V
stroke
LT0
0.00 0.00 0.00 C 3978 3106 M
-27 88 V
stroke
LT0
0.00 0.00 0.00 C 3951 3194 M
-27 88 V
stroke
LT0
0.00 0.00 0.00 C 3924 3282 M
-30 87 V
stroke
LT0
0.00 0.00 0.00 C 3894 3369 M
-31 87 V
stroke
LT0
0.00 0.00 0.00 C 3863 3456 M
-34 87 V
stroke
LT0
0.00 0.00 0.00 C 3829 3543 M
-34 82 V
stroke
LT0
0.00 0.00 0.00 C 3795 3625 M
-1 4 V
stroke
LT0
0.00 0.00 0.00 C 3794 3629 M
-41 86 V
stroke
LT0
0.00 0.00 0.00 C 3753 3715 M
-43 85 V
stroke
LT0
0.00 0.00 0.00 C 3710 3800 M
-46 85 V
stroke
LT0
0.00 0.00 0.00 C 3664 3885 M
-48 84 V
stroke
LT0
0.00 0.00 0.00 C 3616 3969 M
-51 84 V
stroke
LT0
0.00 0.00 0.00 C 3565 4053 M
-27 42 V
stroke
LT0
0.00 0.00 0.00 C 3538 4095 M
-30 41 V
stroke
LT0
0.00 0.00 0.00 C 3508 4136 M
-63 82 V
stroke
LT0
0.00 0.00 0.00 C 3445 4218 M
-65 82 V
4146 2188 M
stroke
LT0
0.00 0.00 0.00 C 4146 2188 M
-15 90 V
stroke
LT0
0.00 0.00 0.00 C 4131 2278 M
-16 90 V
stroke
LT0
0.00 0.00 0.00 C 4115 2368 M
-16 89 V
stroke
LT0
0.00 0.00 0.00 C 4099 2457 M
-17 89 V
stroke
LT0
0.00 0.00 0.00 C 4082 2546 M
-18 90 V
stroke
LT0
0.00 0.00 0.00 C 4064 2636 M
-18 82 V
stroke
LT0
0.00 0.00 0.00 C 4046 2718 M
-1 7 V
stroke
LT0
0.00 0.00 0.00 C 4045 2725 M
-21 88 V
stroke
LT0
0.00 0.00 0.00 C 4024 2813 M
-22 89 V
stroke
LT0
0.00 0.00 0.00 C 4002 2902 M
-23 87 V
stroke
LT0
0.00 0.00 0.00 C 3979 2989 M
-25 88 V
stroke
LT0
0.00 0.00 0.00 C 3954 3077 M
-26 88 V
stroke
LT0
0.00 0.00 0.00 C 3928 3165 M
-28 88 V
stroke
LT0
0.00 0.00 0.00 C 3900 3253 M
-29 87 V
stroke
LT0
0.00 0.00 0.00 C 3871 3340 M
-8 21 V
stroke
LT0
0.00 0.00 0.00 C 3863 3361 M
-25 66 V
stroke
LT0
0.00 0.00 0.00 C 3838 3427 M
-35 87 V
stroke
LT0
0.00 0.00 0.00 C 3803 3514 M
-37 86 V
stroke
LT0
0.00 0.00 0.00 C 3766 3600 M
-41 85 V
stroke
LT0
0.00 0.00 0.00 C 3725 3685 M
-43 85 V
stroke
LT0
0.00 0.00 0.00 C 3682 3770 M
-46 85 V
stroke
LT0
0.00 0.00 0.00 C 3636 3855 M
-48 84 V
stroke
LT0
0.00 0.00 0.00 C 3588 3939 M
-23 39 V
stroke
LT0
0.00 0.00 0.00 C 3565 3978 M
-32 45 V
stroke
LT0
0.00 0.00 0.00 C 3533 4023 M
-60 82 V
stroke
LT0
0.00 0.00 0.00 C 3473 4105 M
-63 82 V
stroke
LT0
0.00 0.00 0.00 C 3410 4187 M
-43 53 V
stroke
LT0
0.00 0.00 0.00 C 3367 4240 M
-28 28 V
4115 2159 M
stroke
LT0
0.00 0.00 0.00 C 4115 2159 M
-15 89 V
stroke
LT0
0.00 0.00 0.00 C 4100 2248 M
-15 90 V
stroke
LT0
0.00 0.00 0.00 C 4085 2338 M
-16 89 V
stroke
LT0
0.00 0.00 0.00 C 4069 2427 M
-17 90 V
stroke
LT0
0.00 0.00 0.00 C 4052 2517 M
-18 89 V
stroke
LT0
0.00 0.00 0.00 C 4034 2606 M
-19 89 V
stroke
LT0
0.00 0.00 0.00 C 4015 2695 M
-19 89 V
stroke
LT0
0.00 0.00 0.00 C 3996 2784 M
-22 88 V
stroke
LT0
0.00 0.00 0.00 C 3974 2872 M
-22 88 V
stroke
LT0
0.00 0.00 0.00 C 3952 2960 M
-25 88 V
stroke
LT0
0.00 0.00 0.00 C 3927 3048 M
-26 88 V
stroke
LT0
0.00 0.00 0.00 C 3901 3136 M
-28 88 V
stroke
LT0
0.00 0.00 0.00 C 3873 3224 M
-30 87 V
stroke
LT0
0.00 0.00 0.00 C 3843 3311 M
-33 87 V
stroke
LT0
0.00 0.00 0.00 C 3810 3398 M
-35 86 V
stroke
LT0
0.00 0.00 0.00 C 3775 3484 M
-37 86 V
stroke
LT0
0.00 0.00 0.00 C 3738 3570 M
-41 86 V
stroke
LT0
0.00 0.00 0.00 C 3697 3656 M
-43 85 V
stroke
LT0
0.00 0.00 0.00 C 3654 3741 M
-46 85 V
stroke
LT0
0.00 0.00 0.00 C 3608 3826 M
-15 25 V
stroke
LT0
0.00 0.00 0.00 C 3593 3851 M
-37 58 V
stroke
LT0
0.00 0.00 0.00 C 3556 3909 M
-58 83 V
stroke
LT0
0.00 0.00 0.00 C 3498 3992 M
-61 83 V
stroke
LT0
0.00 0.00 0.00 C 3437 4075 M
-48 62 V
stroke
LT0
0.00 0.00 0.00 C 3389 4137 M
-18 19 V
stroke
LT0
0.00 0.00 0.00 C 3371 4156 M
-78 80 V
4084 2129 M
stroke
LT0
0.00 0.00 0.00 C 4084 2129 M
-14 89 V
stroke
LT0
0.00 0.00 0.00 C 4070 2218 M
-15 90 V
stroke
LT0
0.00 0.00 0.00 C 4055 2308 M
-16 90 V
stroke
LT0
0.00 0.00 0.00 C 4039 2398 M
-17 89 V
stroke
LT0
0.00 0.00 0.00 C 4022 2487 M
-17 89 V
stroke
LT0
0.00 0.00 0.00 C 4005 2576 M
-19 89 V
stroke
LT0
0.00 0.00 0.00 C 3986 2665 M
-20 89 V
stroke
LT0
0.00 0.00 0.00 C 3966 2754 M
-21 89 V
stroke
LT0
0.00 0.00 0.00 C 3945 2843 M
-22 88 V
stroke
LT0
0.00 0.00 0.00 C 3923 2931 M
-24 88 V
stroke
LT0
0.00 0.00 0.00 C 3899 3019 M
-26 87 V
stroke
LT0
0.00 0.00 0.00 C 3873 3106 M
-28 88 V
stroke
LT0
0.00 0.00 0.00 C 3845 3194 M
-30 87 V
stroke
LT0
0.00 0.00 0.00 C 3815 3281 M
-33 87 V
stroke
LT0
0.00 0.00 0.00 C 3782 3368 M
-35 87 V
stroke
LT0
0.00 0.00 0.00 C 3747 3455 M
-37 86 V
stroke
LT0
0.00 0.00 0.00 C 3710 3541 M
-41 85 V
stroke
LT0
0.00 0.00 0.00 C 3669 3626 M
-43 86 V
stroke
LT0
0.00 0.00 0.00 C 3626 3712 M
-51 83 V
stroke
LT0
0.00 0.00 0.00 C 3575 3795 M
-54 84 V
stroke
LT0
0.00 0.00 0.00 C 3521 3879 M
-58 83 V
stroke
LT0
0.00 0.00 0.00 C 3463 3962 M
-51 68 V
stroke
LT0
0.00 0.00 0.00 C 3412 4030 M
-12 14 V
stroke
LT0
0.00 0.00 0.00 C 3400 4044 M
-75 80 V
stroke
LT0
0.00 0.00 0.00 C 3325 4124 M
-79 79 V
4048 2098 M
stroke
LT0
0.00 0.00 0.00 C 4048 2098 M
-12 90 V
stroke
LT0
0.00 0.00 0.00 C 4036 2188 M
-12 90 V
stroke
LT0
0.00 0.00 0.00 C 4024 2278 M
-4 26 V
stroke
LT0
0.00 0.00 0.00 C 4020 2304 M
-11 64 V
stroke
LT0
0.00 0.00 0.00 C 4009 2368 M
-17 89 V
stroke
LT0
0.00 0.00 0.00 C 3992 2457 M
-17 89 V
stroke
LT0
0.00 0.00 0.00 C 3975 2546 M
-19 90 V
stroke
LT0
0.00 0.00 0.00 C 3956 2636 M
-19 89 V
stroke
LT0
0.00 0.00 0.00 C 3937 2725 M
-21 88 V
stroke
LT0
0.00 0.00 0.00 C 3916 2813 M
-22 89 V
stroke
LT0
0.00 0.00 0.00 C 3894 2902 M
-24 87 V
stroke
LT0
0.00 0.00 0.00 C 3870 2989 M
-26 88 V
stroke
LT0
0.00 0.00 0.00 C 3844 3077 M
-27 88 V
stroke
LT0
0.00 0.00 0.00 C 3817 3165 M
-30 87 V
stroke
LT0
0.00 0.00 0.00 C 3787 3252 M
-33 87 V
stroke
LT0
0.00 0.00 0.00 C 3754 3339 M
-35 86 V
stroke
LT0
0.00 0.00 0.00 C 3719 3425 M
-37 86 V
stroke
LT0
0.00 0.00 0.00 C 3682 3511 M
-15 32 V
stroke
LT0
0.00 0.00 0.00 C 3667 3543 M
-28 54 V
stroke
LT0
0.00 0.00 0.00 C 3639 3597 M
-47 84 V
stroke
LT0
0.00 0.00 0.00 C 3592 3681 M
-51 84 V
stroke
LT0
0.00 0.00 0.00 C 3541 3765 M
-55 83 V
stroke
LT0
0.00 0.00 0.00 C 3486 3848 M
-49 70 V
stroke
LT0
0.00 0.00 0.00 C 3437 3918 M
-11 13 V
stroke
LT0
0.00 0.00 0.00 C 3426 3931 M
-71 81 V
stroke
LT0
0.00 0.00 0.00 C 3355 4012 M
-76 80 V
stroke
LT0
0.00 0.00 0.00 C 3279 4092 M
-20 21 V
stroke
LT0
0.00 0.00 0.00 C 3259 4113 M
-72 56 V
4009 2067 M
stroke
LT0
0.00 0.00 0.00 C 4009 2067 M
-11 90 V
stroke
LT0
0.00 0.00 0.00 C 3998 2157 M
-13 90 V
stroke
LT0
0.00 0.00 0.00 C 3985 2247 M
-13 90 V
stroke
LT0
0.00 0.00 0.00 C 3972 2337 M
-14 90 V
stroke
LT0
0.00 0.00 0.00 C 3958 2427 M
-15 89 V
stroke
LT0
0.00 0.00 0.00 C 3943 2516 M
-16 90 V
stroke
LT0
0.00 0.00 0.00 C 3927 2606 M
-20 89 V
stroke
LT0
0.00 0.00 0.00 C 3907 2695 M
-20 89 V
stroke
LT0
0.00 0.00 0.00 C 3887 2784 M
-22 88 V
stroke
LT0
0.00 0.00 0.00 C 3865 2872 M
-24 87 V
stroke
LT0
0.00 0.00 0.00 C 3841 2959 M
-25 88 V
stroke
LT0
0.00 0.00 0.00 C 3816 3047 M
-28 88 V
stroke
LT0
0.00 0.00 0.00 C 3788 3135 M
-29 87 V
stroke
LT0
0.00 0.00 0.00 C 3759 3222 M
-33 87 V
stroke
LT0
0.00 0.00 0.00 C 3726 3309 M
-37 86 V
stroke
LT0
0.00 0.00 0.00 C 3689 3395 M
-40 86 V
stroke
LT0
0.00 0.00 0.00 C 3649 3481 M
-44 85 V
stroke
LT0
0.00 0.00 0.00 C 3605 3566 M
-48 85 V
stroke
LT0
0.00 0.00 0.00 C 3557 3651 M
-51 83 V
stroke
LT0
0.00 0.00 0.00 C 3506 3734 M
-42 64 V
stroke
LT0
0.00 0.00 0.00 C 3464 3798 M
-15 19 V
stroke
LT0
0.00 0.00 0.00 C 3449 3817 M
-68 82 V
stroke
LT0
0.00 0.00 0.00 C 3381 3899 M
-71 80 V
stroke
LT0
0.00 0.00 0.00 C 3310 3979 M
-29 31 V
stroke
LT0
0.00 0.00 0.00 C 3281 4010 M
-58 48 V
stroke
LT0
0.00 0.00 0.00 C 3223 4058 M
-98 76 V
3970 2036 M
stroke
LT0
0.00 0.00 0.00 C 3970 2036 M
-11 90 V
stroke
LT0
0.00 0.00 0.00 C 3959 2126 M
-12 90 V
stroke
LT0
0.00 0.00 0.00 C 3947 2216 M
-13 90 V
stroke
LT0
0.00 0.00 0.00 C 3934 2306 M
-13 90 V
stroke
LT0
0.00 0.00 0.00 C 3921 2396 M
-15 90 V
stroke
LT0
0.00 0.00 0.00 C 3906 2486 M
-16 89 V
stroke
LT0
0.00 0.00 0.00 C 3890 2575 M
-18 89 V
stroke
LT0
0.00 0.00 0.00 C 3872 2664 M
-19 89 V
stroke
LT0
0.00 0.00 0.00 C 3853 2753 M
-21 89 V
stroke
LT0
0.00 0.00 0.00 C 3832 2842 M
-23 88 V
stroke
LT0
0.00 0.00 0.00 C 3809 2930 M
-25 87 V
stroke
LT0
0.00 0.00 0.00 C 3784 3017 M
-28 88 V
stroke
LT0
0.00 0.00 0.00 C 3756 3105 M
-31 87 V
stroke
LT0
0.00 0.00 0.00 C 3725 3192 M
-33 87 V
stroke
LT0
0.00 0.00 0.00 C 3692 3279 M
-37 86 V
stroke
LT0
0.00 0.00 0.00 C 3655 3365 M
-40 86 V
stroke
LT0
0.00 0.00 0.00 C 3615 3451 M
-44 85 V
stroke
LT0
0.00 0.00 0.00 C 3571 3536 M
-48 84 V
stroke
LT0
0.00 0.00 0.00 C 3523 3620 M
-29 46 V
stroke
LT0
0.00 0.00 0.00 C 3494 3666 M
-26 37 V
stroke
LT0
0.00 0.00 0.00 C 3468 3703 M
-63 82 V
stroke
LT0
0.00 0.00 0.00 C 3405 3785 M
-68 82 V
stroke
LT0
0.00 0.00 0.00 C 3337 3867 M
-33 36 V
stroke
LT0
0.00 0.00 0.00 C 3304 3903 M
-48 43 V
stroke
LT0
0.00 0.00 0.00 C 3256 3946 M
-94 77 V
stroke
LT0
0.00 0.00 0.00 C 3162 4023 M
-20 16 V
stroke
LT0
0.00 0.00 0.00 C 3142 4039 M
-102 57 V
3923 2003 M
stroke
LT0
0.00 0.00 0.00 C 3923 2003 M
-6 91 V
stroke
LT0
0.00 0.00 0.00 C 3917 2094 M
-8 91 V
stroke
LT0
0.00 0.00 0.00 C 3909 2185 M
0 1 V
stroke
LT0
0.00 0.00 0.00 C 3909 2186 M
-12 89 V
stroke
LT0
0.00 0.00 0.00 C 3897 2275 M
-14 90 V
stroke
LT0
0.00 0.00 0.00 C 3883 2365 M
-14 90 V
stroke
LT0
0.00 0.00 0.00 C 3869 2455 M
-16 89 V
stroke
LT0
0.00 0.00 0.00 C 3853 2544 M
-17 89 V
stroke
LT0
0.00 0.00 0.00 C 3836 2633 M
-19 90 V
stroke
LT0
0.00 0.00 0.00 C 3817 2723 M
-21 88 V
stroke
LT0
0.00 0.00 0.00 C 3796 2811 M
-22 89 V
stroke
LT0
0.00 0.00 0.00 C 3774 2900 M
-25 87 V
stroke
LT0
0.00 0.00 0.00 C 3749 2987 M
-28 87 V
stroke
LT0
0.00 0.00 0.00 C 3721 3074 M
-30 88 V
stroke
LT0
0.00 0.00 0.00 C 3691 3162 M
-34 86 V
stroke
LT0
0.00 0.00 0.00 C 3657 3248 M
-36 87 V
stroke
LT0
0.00 0.00 0.00 C 3621 3335 M
-41 85 V
stroke
LT0
0.00 0.00 0.00 C 3580 3420 M
-44 85 V
stroke
LT0
0.00 0.00 0.00 C 3536 3505 M
-6 11 V
stroke
LT0
0.00 0.00 0.00 C 3530 3516 M
-47 73 V
stroke
LT0
0.00 0.00 0.00 C 3483 3589 M
-59 83 V
stroke
LT0
0.00 0.00 0.00 C 3424 3672 M
-63 81 V
stroke
LT0
0.00 0.00 0.00 C 3361 3753 M
-33 39 V
stroke
LT0
0.00 0.00 0.00 C 3328 3792 M
-43 41 V
stroke
LT0
0.00 0.00 0.00 C 3285 3833 M
-89 78 V
stroke
LT0
0.00 0.00 0.00 C 3196 3911 M
-32 27 V
stroke
LT0
0.00 0.00 0.00 C 3164 3938 M
-80 48 V
stroke
LT0
0.00 0.00 0.00 C 3084 3986 M
-68 37 V
stroke
LT0
0.00 0.00 0.00 C 3016 4023 M
-87 30 V
3873 1970 M
stroke
LT0
0.00 0.00 0.00 C 3873 1970 M
-6 91 V
stroke
LT0
0.00 0.00 0.00 C 3867 2061 M
-8 91 V
stroke
LT0
0.00 0.00 0.00 C 3859 2152 M
-8 91 V
stroke
LT0
0.00 0.00 0.00 C 3851 2243 M
-9 90 V
stroke
LT0
0.00 0.00 0.00 C 3842 2333 M
-12 91 V
stroke
LT0
0.00 0.00 0.00 C 3830 2424 M
-3 25 V
stroke
LT0
0.00 0.00 0.00 C 3827 2449 M
-11 64 V
stroke
LT0
0.00 0.00 0.00 C 3816 2513 M
-17 90 V
stroke
LT0
0.00 0.00 0.00 C 3799 2603 M
-18 89 V
stroke
LT0
0.00 0.00 0.00 C 3781 2692 M
-21 89 V
stroke
LT0
0.00 0.00 0.00 C 3760 2781 M
-22 88 V
stroke
LT0
0.00 0.00 0.00 C 3738 2869 M
-25 87 V
stroke
LT0
0.00 0.00 0.00 C 3713 2956 M
-27 88 V
stroke
LT0
0.00 0.00 0.00 C 3686 3044 M
-30 87 V
stroke
LT0
0.00 0.00 0.00 C 3656 3131 M
-33 87 V
stroke
LT0
0.00 0.00 0.00 C 3623 3218 M
-37 86 V
stroke
LT0
0.00 0.00 0.00 C 3586 3304 M
-7 17 V
stroke
LT0
0.00 0.00 0.00 C 3579 3321 M
-36 68 V
stroke
LT0
0.00 0.00 0.00 C 3543 3389 M
-49 85 V
stroke
LT0
0.00 0.00 0.00 C 3494 3474 M
-54 83 V
stroke
LT0
0.00 0.00 0.00 C 3440 3557 M
-59 83 V
stroke
LT0
0.00 0.00 0.00 C 3381 3640 M
-26 32 V
stroke
LT0
0.00 0.00 0.00 C 3355 3672 M
-46 48 V
stroke
LT0
0.00 0.00 0.00 C 3309 3720 M
-83 79 V
stroke
LT0
0.00 0.00 0.00 C 3226 3799 M
-40 35 V
stroke
LT0
0.00 0.00 0.00 C 3186 3834 M
-63 41 V
stroke
LT0
0.00 0.00 0.00 C 3123 3875 M
-86 51 V
stroke
LT0
0.00 0.00 0.00 C 3037 3926 M
-51 19 V
stroke
LT0
0.00 0.00 0.00 C 2986 3945 M
-86 30 V
stroke
LT0
0.00 0.00 0.00 C 2900 3975 M
-128 19 V
stroke
LT0
0.00 0.00 0.00 C 2772 3994 M
-59 0 V
3816 1936 M
stroke
LT0
0.00 0.00 0.00 C 3816 1936 M
1 90 V
stroke
LT0
0.00 0.00 0.00 C 3817 2026 M
0 2 V
stroke
LT0
0.00 0.00 0.00 C 3817 2028 M
-7 91 V
stroke
LT0
0.00 0.00 0.00 C 3810 2119 M
-8 91 V
stroke
LT0
0.00 0.00 0.00 C 3802 2210 M
-9 91 V
stroke
LT0
0.00 0.00 0.00 C 3793 2301 M
-11 90 V
stroke
LT0
0.00 0.00 0.00 C 3782 2391 M
-12 90 V
stroke
LT0
0.00 0.00 0.00 C 3770 2481 M
-14 90 V
stroke
LT0
0.00 0.00 0.00 C 3756 2571 M
-16 89 V
stroke
LT0
0.00 0.00 0.00 C 3740 2660 M
-18 89 V
stroke
LT0
0.00 0.00 0.00 C 3722 2749 M
-21 89 V
stroke
LT0
0.00 0.00 0.00 C 3701 2838 M
-24 88 V
stroke
LT0
0.00 0.00 0.00 C 3677 2926 M
-28 87 V
stroke
LT0
0.00 0.00 0.00 C 3649 3013 M
-31 87 V
stroke
LT0
0.00 0.00 0.00 C 3618 3100 M
-35 87 V
stroke
LT0
0.00 0.00 0.00 C 3583 3187 M
-39 85 V
stroke
LT0
0.00 0.00 0.00 C 3544 3272 M
-44 85 V
stroke
LT0
0.00 0.00 0.00 C 3500 3357 M
-49 85 V
stroke
LT0
0.00 0.00 0.00 C 3451 3442 M
-55 83 V
stroke
LT0
0.00 0.00 0.00 C 3396 3525 M
-10 15 V
stroke
LT0
0.00 0.00 0.00 C 3386 3540 M
-58 66 V
stroke
LT0
0.00 0.00 0.00 C 3328 3606 M
-77 80 V
stroke
LT0
0.00 0.00 0.00 C 3251 3686 M
-41 39 V
stroke
LT0
0.00 0.00 0.00 C 3210 3725 M
-54 38 V
stroke
LT0
0.00 0.00 0.00 C 3156 3763 M
-98 63 V
stroke
LT0
0.00 0.00 0.00 C 3058 3826 M
-24 10 V
stroke
LT0
0.00 0.00 0.00 C 3034 3836 M
-114 43 V
stroke
LT0
0.00 0.00 0.00 C 2920 3879 M
-93 16 V
stroke
LT0
0.00 0.00 0.00 C 2827 3895 M
-35 5 V
stroke
LT0
0.00 0.00 0.00 C 2792 3900 M
-124 0 V
stroke
LT0
0.00 0.00 0.00 C 2668 3900 M
-119 -14 V
stroke
LT0
0.00 0.00 0.00 C 2549 3886 M
-117 -21 V
stroke
LT0
0.00 0.00 0.00 C 2432 3865 M
-117 -22 V
stroke
LT0
0.00 0.00 0.00 C 2315 3843 M
-118 -17 V
stroke
LT0
0.00 0.00 0.00 C 2197 3826 M
-123 -2 V
stroke
LT0
0.00 0.00 0.00 C 2074 3824 M
-128 19 V
stroke
LT0
0.00 0.00 0.00 C 1946 3843 M
-7 3 V
3748 1900 M
stroke
LT0
0.00 0.00 0.00 C 3748 1900 M
1 93 V
stroke
LT0
0.00 0.00 0.00 C 3749 1993 M
1 92 V
stroke
LT0
0.00 0.00 0.00 C 3750 2085 M
-2 92 V
stroke
LT0
0.00 0.00 0.00 C 3748 2177 M
-2 73 V
stroke
LT0
0.00 0.00 0.00 C 3746 2250 M
-2 18 V
stroke
LT0
0.00 0.00 0.00 C 3744 2268 M
-10 90 V
stroke
LT0
0.00 0.00 0.00 C 3734 2358 M
-11 91 V
stroke
LT0
0.00 0.00 0.00 C 3723 2449 M
-14 89 V
stroke
LT0
0.00 0.00 0.00 C 3709 2538 M
-15 90 V
stroke
LT0
0.00 0.00 0.00 C 3694 2628 M
-18 89 V
stroke
LT0
0.00 0.00 0.00 C 3676 2717 M
-21 89 V
stroke
LT0
0.00 0.00 0.00 C 3655 2806 M
-23 88 V
stroke
LT0
0.00 0.00 0.00 C 3632 2894 M
-27 87 V
stroke
LT0
0.00 0.00 0.00 C 3605 2981 M
-30 87 V
stroke
LT0
0.00 0.00 0.00 C 3575 3068 M
-35 87 V
stroke
LT0
0.00 0.00 0.00 C 3540 3155 M
-39 86 V
stroke
LT0
0.00 0.00 0.00 C 3501 3241 M
-44 85 V
stroke
LT0
0.00 0.00 0.00 C 3457 3326 M
-33 57 V
stroke
LT0
0.00 0.00 0.00 C 3424 3383 M
-19 26 V
stroke
LT0
0.00 0.00 0.00 C 3405 3409 M
-63 82 V
stroke
LT0
0.00 0.00 0.00 C 3342 3491 M
-70 81 V
stroke
LT0
0.00 0.00 0.00 C 3272 3572 M
-36 37 V
stroke
LT0
0.00 0.00 0.00 C 3236 3609 M
-53 41 V
stroke
LT0
0.00 0.00 0.00 C 3183 3650 M
-102 72 V
stroke
LT0
0.00 0.00 0.00 C 3081 3722 M
-7 3 V
stroke
LT0
0.00 0.00 0.00 C 3074 3725 M
-133 55 V
stroke
LT0
0.00 0.00 0.00 C 2941 3780 M
-44 9 V
stroke
LT0
0.00 0.00 0.00 C 2897 3789 M
-85 15 V
stroke
LT0
0.00 0.00 0.00 C 2812 3804 M
-124 2 V
stroke
LT0
0.00 0.00 0.00 C 2688 3806 M
-119 -14 V
stroke
LT0
0.00 0.00 0.00 C 2569 3792 M
-117 -22 V
stroke
LT0
0.00 0.00 0.00 C 2452 3770 M
-117 -23 V
stroke
LT0
0.00 0.00 0.00 C 2335 3747 M
-118 -17 V
stroke
LT0
0.00 0.00 0.00 C 2217 3730 M
-122 -4 V
stroke
LT0
0.00 0.00 0.00 C 2095 3726 M
-89 13 V
stroke
LT0
0.00 0.00 0.00 C 2006 3739 M
-40 6 V
stroke
LT0
0.00 0.00 0.00 C 1966 3745 M
-137 49 V
stroke
LT0
0.00 0.00 0.00 C 1829 3794 M
-12 7 V
3661 1862 M
stroke
LT0
0.00 0.00 0.00 C 3661 1862 M
14 94 V
stroke
LT0
0.00 0.00 0.00 C 3675 1956 M
7 60 V
stroke
LT0
0.00 0.00 0.00 C 3682 2016 M
1 33 V
stroke
LT0
0.00 0.00 0.00 C 3683 2049 M
-1 92 V
stroke
LT0
0.00 0.00 0.00 C 3682 2141 M
-2 92 V
stroke
LT0
0.00 0.00 0.00 C 3680 2233 M
-4 91 V
stroke
LT0
0.00 0.00 0.00 C 3676 2324 M
-6 91 V
stroke
LT0
0.00 0.00 0.00 C 3670 2415 M
-9 91 V
stroke
LT0
0.00 0.00 0.00 C 3661 2506 M
-4 34 V
stroke
LT0
0.00 0.00 0.00 C 3657 2540 M
-10 56 V
stroke
LT0
0.00 0.00 0.00 C 3647 2596 M
-17 89 V
stroke
LT0
0.00 0.00 0.00 C 3630 2685 M
-20 89 V
stroke
LT0
0.00 0.00 0.00 C 3610 2774 M
-23 88 V
stroke
LT0
0.00 0.00 0.00 C 3587 2862 M
-26 87 V
stroke
LT0
0.00 0.00 0.00 C 3561 2949 M
-30 87 V
stroke
LT0
0.00 0.00 0.00 C 3531 3036 M
-34 87 V
stroke
LT0
0.00 0.00 0.00 C 3497 3123 M
-17 37 V
stroke
LT0
0.00 0.00 0.00 C 3480 3160 M
-25 48 V
stroke
LT0
0.00 0.00 0.00 C 3455 3208 M
-49 85 V
stroke
LT0
0.00 0.00 0.00 C 3406 3293 M
-56 83 V
stroke
LT0
0.00 0.00 0.00 C 3350 3376 M
-64 82 V
stroke
LT0
0.00 0.00 0.00 C 3286 3458 M
-21 23 V
stroke
LT0
0.00 0.00 0.00 C 3265 3481 M
-62 55 V
stroke
LT0
0.00 0.00 0.00 C 3203 3536 M
-98 77 V
stroke
LT0
0.00 0.00 0.00 C 3105 3613 M
-142 65 V
stroke
LT0
0.00 0.00 0.00 C 2963 3678 M
-10 3 V
stroke
LT0
0.00 0.00 0.00 C 2953 3681 M
-120 25 V
stroke
LT0
0.00 0.00 0.00 C 2833 3706 M
-124 3 V
stroke
LT0
0.00 0.00 0.00 C 2709 3709 M
-120 -13 V
stroke
LT0
0.00 0.00 0.00 C 2589 3696 M
-117 -22 V
stroke
LT0
0.00 0.00 0.00 C 2472 3674 M
-116 -24 V
stroke
LT0
0.00 0.00 0.00 C 2356 3650 M
-118 -19 V
stroke
LT0
0.00 0.00 0.00 C 2238 3631 M
-122 -5 V
stroke
LT0
0.00 0.00 0.00 C 2116 3626 M
-45 7 V
stroke
LT0
0.00 0.00 0.00 C 2071 3633 M
-83 11 V
stroke
LT0
0.00 0.00 0.00 C 1988 3644 M
-129 47 V
stroke
LT0
0.00 0.00 0.00 C 1859 3691 M
-9 3 V
stroke
LT0
0.00 0.00 0.00 C 1850 3694 M
-115 69 V
3540 1817 M
stroke
LT0
0.00 0.00 0.00 C 3540 1817 M
37 98 V
stroke
LT0
0.00 0.00 0.00 C 3577 1915 M
2 6 V
stroke
LT0
0.00 0.00 0.00 C 3579 1921 M
13 89 V
stroke
LT0
0.00 0.00 0.00 C 3592 2010 M
11 93 V
stroke
LT0
0.00 0.00 0.00 C 3603 2103 M
8 94 V
stroke
LT0
0.00 0.00 0.00 C 3611 2197 M
3 36 V
stroke
LT0
0.00 0.00 0.00 C 3614 2233 M
-3 56 V
stroke
LT0
0.00 0.00 0.00 C 3611 2289 M
-5 91 V
stroke
LT0
0.00 0.00 0.00 C 3606 2380 M
-8 91 V
stroke
LT0
0.00 0.00 0.00 C 3598 2471 M
-10 90 V
stroke
LT0
0.00 0.00 0.00 C 3588 2561 M
-14 90 V
stroke
LT0
0.00 0.00 0.00 C 3574 2651 M
-17 89 V
stroke
LT0
0.00 0.00 0.00 C 3557 2740 M
-22 89 V
stroke
LT0
0.00 0.00 0.00 C 3535 2829 M
-25 88 V
stroke
LT0
0.00 0.00 0.00 C 3510 2917 M
-31 86 V
stroke
LT0
0.00 0.00 0.00 C 3479 3003 M
-37 86 V
stroke
LT0
0.00 0.00 0.00 C 3442 3089 M
-42 86 V
stroke
LT0
0.00 0.00 0.00 C 3400 3175 M
-50 84 V
stroke
LT0
0.00 0.00 0.00 C 3350 3259 M
-49 73 V
stroke
LT0
0.00 0.00 0.00 C 3301 3332 M
-9 10 V
stroke
LT0
0.00 0.00 0.00 C 3292 3342 M
-77 79 V
stroke
LT0
0.00 0.00 0.00 C 3215 3421 M
-84 74 V
stroke
LT0
0.00 0.00 0.00 C 3131 3495 M
-7 4 V
stroke
LT0
0.00 0.00 0.00 C 3124 3499 M
-134 71 V
stroke
LT0
0.00 0.00 0.00 C 2990 3570 M
-4 1 V
stroke
LT0
0.00 0.00 0.00 C 2986 3571 M
-132 33 V
stroke
LT0
0.00 0.00 0.00 C 2854 3604 M
-124 5 V
stroke
LT0
0.00 0.00 0.00 C 2730 3609 M
-120 -12 V
stroke
LT0
0.00 0.00 0.00 C 2610 3597 M
-117 -23 V
stroke
LT0
0.00 0.00 0.00 C 2493 3574 M
-115 -25 V
stroke
LT0
0.00 0.00 0.00 C 2378 3549 M
-118 -21 V
stroke
LT0
0.00 0.00 0.00 C 2260 3528 M
-121 -6 V
stroke
LT0
0.00 0.00 0.00 C 2139 3522 M
-13 2 V
stroke
LT0
0.00 0.00 0.00 C 2126 3524 M
-115 16 V
stroke
LT0
0.00 0.00 0.00 C 2011 3540 M
-110 41 V
stroke
LT0
0.00 0.00 0.00 C 1901 3581 M
-28 10 V
stroke
LT0
0.00 0.00 0.00 C 1873 3591 M
-101 61 V
stroke
LT0
0.00 0.00 0.00 C 1772 3652 M
-48 29 V
stroke
LT0
0.00 0.00 0.00 C 1724 3681 M
-55 47 V
3211 1740 M
stroke
LT0
0.00 0.00 0.00 C 3211 1740 M
40 10 V
stroke
LT0
0.00 0.00 0.00 C 3251 1750 M
107 59 V
stroke
LT0
0.00 0.00 0.00 C 3358 1809 M
44 54 V
stroke
LT0
0.00 0.00 0.00 C 3402 1863 M
47 58 V
stroke
LT0
0.00 0.00 0.00 C 3449 1921 M
16 43 V
stroke
LT0
0.00 0.00 0.00 C 3465 1964 M
32 98 V
stroke
LT0
0.00 0.00 0.00 C 3497 2062 M
18 59 V
stroke
LT0
0.00 0.00 0.00 C 3515 2121 M
4 36 V
stroke
LT0
0.00 0.00 0.00 C 3519 2157 M
7 93 V
stroke
LT0
0.00 0.00 0.00 C 3526 2250 M
4 93 V
stroke
LT0
0.00 0.00 0.00 C 3530 2343 M
1 92 V
stroke
LT0
0.00 0.00 0.00 C 3531 2435 M
-3 74 V
stroke
LT0
0.00 0.00 0.00 C 3528 2509 M
-2 17 V
stroke
LT0
0.00 0.00 0.00 C 3526 2526 M
-13 90 V
stroke
LT0
0.00 0.00 0.00 C 3513 2616 M
-16 90 V
stroke
LT0
0.00 0.00 0.00 C 3497 2706 M
-20 89 V
stroke
LT0
0.00 0.00 0.00 C 3477 2795 M
-25 88 V
stroke
LT0
0.00 0.00 0.00 C 3452 2883 M
-30 86 V
stroke
LT0
0.00 0.00 0.00 C 3422 2969 M
-36 86 V
stroke
LT0
0.00 0.00 0.00 C 3386 3055 M
-33 66 V
stroke
LT0
0.00 0.00 0.00 C 3353 3121 M
-11 19 V
stroke
LT0
0.00 0.00 0.00 C 3342 3140 M
-57 83 V
stroke
LT0
0.00 0.00 0.00 C 3285 3223 M
-68 82 V
stroke
LT0
0.00 0.00 0.00 C 3217 3305 M
-55 56 V
stroke
LT0
0.00 0.00 0.00 C 3162 3361 M
-30 22 V
stroke
LT0
0.00 0.00 0.00 C 3132 3383 M
-119 73 V
stroke
LT0
0.00 0.00 0.00 C 3013 3456 M
-1 1 V
stroke
LT0
0.00 0.00 0.00 C 3012 3457 M
-135 39 V
stroke
LT0
0.00 0.00 0.00 C 2877 3496 M
-125 9 V
stroke
LT0
0.00 0.00 0.00 C 2752 3505 M
-21 -2 V
stroke
LT0
0.00 0.00 0.00 C 2731 3503 M
-99 -9 V
stroke
LT0
0.00 0.00 0.00 C 2632 3494 M
-116 -24 V
stroke
LT0
0.00 0.00 0.00 C 2516 3470 M
-18 -4 V
stroke
LT0
0.00 0.00 0.00 C 2498 3466 M
-98 -23 V
stroke
LT0
0.00 0.00 0.00 C 2400 3443 M
-116 -23 V
stroke
LT0
0.00 0.00 0.00 C 2284 3420 M
-121 -7 V
stroke
LT0
0.00 0.00 0.00 C 2163 3413 M
-128 17 V
stroke
LT0
0.00 0.00 0.00 C 2035 3430 M
-100 39 V
stroke
LT0
0.00 0.00 0.00 C 1935 3469 M
-39 14 V
stroke
LT0
0.00 0.00 0.00 C 1896 3483 M
-90 57 V
stroke
LT0
0.00 0.00 0.00 C 1806 3540 M
-60 37 V
stroke
LT0
0.00 0.00 0.00 C 1746 3577 M
-44 39 V
stroke
LT0
0.00 0.00 0.00 C 1702 3616 M
-91 78 V
2590 1616 M
stroke
LT0
0.00 0.00 0.00 C 2590 1616 M
45 99 V
stroke
LT0
0.00 0.00 0.00 C 2635 1715 M
2 6 V
stroke
LT0
0.00 0.00 0.00 C 2637 1721 M
96 96 V
stroke
LT0
0.00 0.00 0.00 C 2733 1817 M
22 10 V
stroke
LT0
0.00 0.00 0.00 C 2755 1827 M
87 37 V
stroke
LT0
0.00 0.00 0.00 C 2842 1864 M
117 21 V
stroke
LT0
0.00 0.00 0.00 C 2959 1885 M
119 15 V
stroke
LT0
0.00 0.00 0.00 C 3078 1900 M
115 29 V
stroke
LT0
0.00 0.00 0.00 C 3193 1929 M
105 63 V
stroke
LT0
0.00 0.00 0.00 C 3298 1992 M
11 15 V
stroke
LT0
0.00 0.00 0.00 C 3309 2007 M
73 103 V
stroke
LT0
0.00 0.00 0.00 C 3382 2110 M
4 7 V
stroke
LT0
0.00 0.00 0.00 C 3386 2117 M
25 90 V
stroke
LT0
0.00 0.00 0.00 C 3411 2207 M
22 96 V
stroke
LT0
0.00 0.00 0.00 C 3433 2303 M
9 50 V
stroke
LT0
0.00 0.00 0.00 C 3442 2353 M
1 43 V
stroke
LT0
0.00 0.00 0.00 C 3443 2396 M
-2 92 V
stroke
LT0
0.00 0.00 0.00 C 3441 2488 M
-6 91 V
stroke
LT0
0.00 0.00 0.00 C 3435 2579 M
-11 90 V
stroke
LT0
0.00 0.00 0.00 C 3424 2669 M
-17 90 V
stroke
LT0
0.00 0.00 0.00 C 3407 2759 M
-23 88 V
stroke
LT0
0.00 0.00 0.00 C 3384 2847 M
-31 87 V
stroke
LT0
0.00 0.00 0.00 C 3353 2934 M
-38 85 V
stroke
LT0
0.00 0.00 0.00 C 3315 3019 M
-48 85 V
stroke
LT0
0.00 0.00 0.00 C 3267 3104 M
-57 82 V
stroke
LT0
0.00 0.00 0.00 C 3210 3186 M
-8 11 V
stroke
LT0
0.00 0.00 0.00 C 3202 3197 M
-76 68 V
stroke
LT0
0.00 0.00 0.00 C 3126 3265 M
-85 62 V
stroke
LT0
0.00 0.00 0.00 C 3041 3327 M
-27 12 V
stroke
LT0
0.00 0.00 0.00 C 3014 3339 M
-111 40 V
stroke
LT0
0.00 0.00 0.00 C 2903 3379 M
-125 15 V
stroke
LT0
0.00 0.00 0.00 C 2778 3394 M
-1 0 V
stroke
LT0
0.00 0.00 0.00 C 2777 3394 M
-120 -11 V
stroke
LT0
0.00 0.00 0.00 C 2657 3383 M
-117 -24 V
stroke
LT0
0.00 0.00 0.00 C 2540 3359 M
-34 -8 V
stroke
LT0
0.00 0.00 0.00 C 2506 3351 M
-80 -22 V
stroke
LT0
0.00 0.00 0.00 C 2426 3329 M
-116 -25 V
stroke
LT0
0.00 0.00 0.00 C 2310 3304 M
-121 -9 V
stroke
LT0
0.00 0.00 0.00 C 2189 3295 M
-17 3 V
stroke
LT0
0.00 0.00 0.00 C 2172 3298 M
-111 15 V
stroke
LT0
0.00 0.00 0.00 C 2061 3313 M
-103 43 V
stroke
LT0
0.00 0.00 0.00 C 1958 3356 M
-37 14 V
stroke
LT0
0.00 0.00 0.00 C 1921 3370 M
-87 58 V
stroke
LT0
0.00 0.00 0.00 C 1834 3428 M
-64 41 V
stroke
LT0
0.00 0.00 0.00 C 1770 3469 M
-38 35 V
stroke
LT0
0.00 0.00 0.00 C 1732 3504 M
-89 78 V
stroke
LT0
0.00 0.00 0.00 C 1643 3582 M
-37 32 V
stroke
LT0
0.00 0.00 0.00 C 1606 3614 M
-43 47 V
2468 1572 M
stroke
LT0
0.00 0.00 0.00 C 2468 1572 M
16 95 V
stroke
LT0
0.00 0.00 0.00 C 2484 1667 M
15 83 V
stroke
LT0
0.00 0.00 0.00 C 2499 1750 M
5 12 V
stroke
LT0
0.00 0.00 0.00 C 2504 1762 M
43 99 V
stroke
LT0
0.00 0.00 0.00 C 2547 1861 M
26 61 V
stroke
LT0
0.00 0.00 0.00 C 2573 1922 M
43 42 V
stroke
LT0
0.00 0.00 0.00 C 2616 1964 M
53 52 V
stroke
LT0
0.00 0.00 0.00 C 2669 2016 M
111 43 V
stroke
LT0
0.00 0.00 0.00 C 2780 2059 M
118 17 V
stroke
LT0
0.00 0.00 0.00 C 2898 2076 M
119 13 V
stroke
LT0
0.00 0.00 0.00 C 3017 2089 M
114 31 V
stroke
LT0
0.00 0.00 0.00 C 3131 2120 M
34 23 V
stroke
LT0
0.00 0.00 0.00 C 3165 2143 M
69 50 V
stroke
LT0
0.00 0.00 0.00 C 3234 2193 M
35 59 V
stroke
LT0
0.00 0.00 0.00 C 3269 2252 M
46 90 V
stroke
LT0
0.00 0.00 0.00 C 3315 2342 M
2 9 V
stroke
LT0
0.00 0.00 0.00 C 3317 2351 M
13 94 V
stroke
LT0
0.00 0.00 0.00 C 3330 2445 M
6 93 V
stroke
LT0
0.00 0.00 0.00 C 3336 2538 M
-1 92 V
stroke
LT0
0.00 0.00 0.00 C 3335 2630 M
-10 91 V
stroke
LT0
0.00 0.00 0.00 C 3325 2721 M
-20 88 V
stroke
LT0
0.00 0.00 0.00 C 3305 2809 M
-31 88 V
stroke
LT0
0.00 0.00 0.00 C 3274 2897 M
-42 84 V
stroke
LT0
0.00 0.00 0.00 C 3232 2981 M
-56 83 V
stroke
LT0
0.00 0.00 0.00 C 3176 3064 M
-70 81 V
stroke
LT0
0.00 0.00 0.00 C 3106 3145 M
-27 24 V
stroke
LT0
0.00 0.00 0.00 C 3079 3169 M
-85 50 V
stroke
LT0
0.00 0.00 0.00 C 2994 3219 M
-60 27 V
stroke
LT0
0.00 0.00 0.00 C 2934 3246 M
-129 23 V
stroke
LT0
0.00 0.00 0.00 C 2805 3269 M
-121 -7 V
stroke
LT0
0.00 0.00 0.00 C 2684 3262 M
-115 -27 V
stroke
LT0
0.00 0.00 0.00 C 2569 3235 M
-114 -34 V
stroke
LT0
0.00 0.00 0.00 C 2455 3201 M
-115 -30 V
stroke
LT0
0.00 0.00 0.00 C 2340 3171 M
-120 -10 V
stroke
LT0
0.00 0.00 0.00 C 2220 3161 M
-78 15 V
stroke
LT0
0.00 0.00 0.00 C 2142 3176 M
-52 8 V
stroke
LT0
0.00 0.00 0.00 C 2090 3184 M
-123 56 V
stroke
LT0
0.00 0.00 0.00 C 1967 3240 M
-18 8 V
stroke
LT0
0.00 0.00 0.00 C 1949 3248 M
-94 66 V
stroke
LT0
0.00 0.00 0.00 C 1855 3314 M
-60 40 V
stroke
LT0
0.00 0.00 0.00 C 1795 3354 M
-38 37 V
stroke
LT0
0.00 0.00 0.00 C 1757 3391 M
-85 78 V
stroke
LT0
0.00 0.00 0.00 C 1672 3469 M
-43 38 V
stroke
LT0
0.00 0.00 0.00 C 1629 3507 M
-37 42 V
stroke
LT0
0.00 0.00 0.00 C 1592 3549 M
-72 80 V
2383 1533 M
stroke
LT0
0.00 0.00 0.00 C 2383 1533 M
2 93 V
stroke
LT0
0.00 0.00 0.00 C 2385 1626 M
1 60 V
stroke
LT0
0.00 0.00 0.00 C 2386 1686 M
6 33 V
stroke
LT0
0.00 0.00 0.00 C 2392 1719 M
14 94 V
stroke
LT0
0.00 0.00 0.00 C 2406 1813 M
14 95 V
stroke
LT0
0.00 0.00 0.00 C 2420 1908 M
9 62 V
stroke
LT0
0.00 0.00 0.00 C 2429 1970 M
13 33 V
stroke
LT0
0.00 0.00 0.00 C 2442 2003 M
38 98 V
stroke
LT0
0.00 0.00 0.00 C 2480 2101 M
20 53 V
stroke
LT0
0.00 0.00 0.00 C 2500 2154 M
51 50 V
stroke
LT0
0.00 0.00 0.00 C 2551 2204 M
45 46 V
stroke
LT0
0.00 0.00 0.00 C 2596 2250 M
113 35 V
stroke
LT0
0.00 0.00 0.00 C 2709 2285 M
121 9 V
stroke
LT0
0.00 0.00 0.00 C 2830 2294 M
119 10 V
stroke
LT0
0.00 0.00 0.00 C 2949 2304 M
113 36 V
stroke
LT0
0.00 0.00 0.00 C 3062 2340 M
51 46 V
stroke
LT0
0.00 0.00 0.00 C 3113 2386 M
45 48 V
stroke
LT0
0.00 0.00 0.00 C 3158 2434 M
19 54 V
stroke
LT0
0.00 0.00 0.00 C 3177 2488 M
21 95 V
stroke
LT0
0.00 0.00 0.00 C 3198 2583 M
5 93 V
stroke
LT0
0.00 0.00 0.00 C 3203 2676 M
-6 56 V
stroke
LT0
0.00 0.00 0.00 C 3197 2732 M
-7 34 V
stroke
LT0
0.00 0.00 0.00 C 3190 2766 M
-27 88 V
stroke
LT0
0.00 0.00 0.00 C 3163 2854 M
-2 4 V
stroke
LT0
0.00 0.00 0.00 C 3161 2858 M
-47 79 V
stroke
LT0
0.00 0.00 0.00 C 3114 2937 M
-70 81 V
stroke
LT0
0.00 0.00 0.00 C 3044 3018 M
-67 54 V
stroke
LT0
0.00 0.00 0.00 C 2977 3072 M
-48 20 V
stroke
LT0
0.00 0.00 0.00 C 2929 3092 M
-88 23 V
stroke
LT0
0.00 0.00 0.00 C 2841 3115 M
-121 -3 V
stroke
LT0
0.00 0.00 0.00 C 2720 3112 M
-115 -29 V
stroke
LT0
0.00 0.00 0.00 C 2605 3083 M
-112 -41 V
stroke
LT0
0.00 0.00 0.00 C 2493 3042 M
-113 -35 V
stroke
LT0
0.00 0.00 0.00 C 2380 3007 M
-120 -10 V
stroke
LT0
0.00 0.00 0.00 C 2260 2997 M
-133 33 V
stroke
LT0
0.00 0.00 0.00 C 2127 3030 M
-33 21 V
stroke
LT0
0.00 0.00 0.00 C 2094 3051 M
-113 59 V
stroke
LT0
0.00 0.00 0.00 C 1981 3110 M
-14 13 V
stroke
LT0
0.00 0.00 0.00 C 1967 3123 M
-98 76 V
stroke
LT0
0.00 0.00 0.00 C 1869 3199 M
-46 33 V
stroke
LT0
0.00 0.00 0.00 C 1823 3232 M
-44 45 V
stroke
LT0
0.00 0.00 0.00 C 1779 3277 M
-82 79 V
stroke
LT0
0.00 0.00 0.00 C 1697 3356 M
-44 40 V
stroke
LT0
0.00 0.00 0.00 C 1653 3396 M
-34 40 V
stroke
LT0
0.00 0.00 0.00 C 1619 3436 M
-70 81 V
stroke
LT0
0.00 0.00 0.00 C 1549 3517 M
-72 80 V
2720 2654 M
stroke
LT0
0.00 0.00 0.00 C 2720 2654 M
-21 9 V
stroke
LT0
0.00 0.00 0.00 C 2699 2663 M
-42 85 V
stroke
LT0
0.00 0.00 0.00 C 2657 2748 M
26 36 V
stroke
LT0
0.00 0.00 0.00 C 2683 2784 M
101 77 V
stroke
LT0
0.00 0.00 0.00 C 2784 2861 M
141 -67 V
stroke
LT0
0.00 0.00 0.00 C 2925 2794 M
3 -2 V
stroke
LT0
0.00 0.00 0.00 C 2928 2792 M
3 -17 V
stroke
LT0
0.00 0.00 0.00 C 2931 2775 M
-1 -75 V
stroke
LT0
0.00 0.00 0.00 C 2930 2700 M
-82 -67 V
stroke
LT0
0.00 0.00 0.00 C 2848 2633 M
-128 21 V
2310 1497 M
stroke
LT0
0.00 0.00 0.00 C 2310 1497 M
2 92 V
stroke
LT0
0.00 0.00 0.00 C 2312 1589 M
1 92 V
stroke
LT0
0.00 0.00 0.00 C 2313 1681 M
1 93 V
stroke
LT0
0.00 0.00 0.00 C 2314 1774 M
0 92 V
stroke
LT0
0.00 0.00 0.00 C 2314 1866 M
-1 54 V
stroke
LT0
0.00 0.00 0.00 C 2313 1920 M
4 38 V
stroke
LT0
0.00 0.00 0.00 C 2317 1958 M
10 94 V
stroke
LT0
0.00 0.00 0.00 C 2327 2052 M
7 93 V
stroke
LT0
0.00 0.00 0.00 C 2334 2145 M
5 93 V
stroke
LT0
0.00 0.00 0.00 C 2339 2238 M
1 21 V
stroke
LT0
0.00 0.00 0.00 C 2340 2259 M
17 74 V
stroke
LT0
0.00 0.00 0.00 C 2357 2333 M
15 94 V
stroke
LT0
0.00 0.00 0.00 C 2372 2427 M
4 93 V
stroke
LT0
0.00 0.00 0.00 C 2376 2520 M
-13 90 V
stroke
LT0
0.00 0.00 0.00 C 2363 2610 M
-41 85 V
stroke
LT0
0.00 0.00 0.00 C 2322 2695 M
-79 80 V
stroke
LT0
0.00 0.00 0.00 C 2243 2775 M
-60 33 V
stroke
LT0
0.00 0.00 0.00 C 2183 2808 M
-39 43 V
stroke
LT0
0.00 0.00 0.00 C 2144 2851 M
-96 77 V
stroke
LT0
0.00 0.00 0.00 C 2048 2928 M
-27 16 V
stroke
LT0
0.00 0.00 0.00 C 2021 2944 M
-60 61 V
stroke
LT0
0.00 0.00 0.00 C 1961 3005 M
-89 78 V
stroke
LT0
0.00 0.00 0.00 C 1872 3083 M
-17 14 V
stroke
LT0
0.00 0.00 0.00 C 1855 3097 M
-60 66 V
stroke
LT0
0.00 0.00 0.00 C 1795 3163 M
-78 79 V
stroke
LT0
0.00 0.00 0.00 C 1717 3242 M
-37 36 V
stroke
LT0
0.00 0.00 0.00 C 1680 3278 M
-36 45 V
stroke
LT0
0.00 0.00 0.00 C 1644 3323 M
-69 81 V
stroke
LT0
0.00 0.00 0.00 C 1575 3404 M
-70 81 V
stroke
LT0
0.00 0.00 0.00 C 1505 3485 M
-11 12 V
stroke
LT0
0.00 0.00 0.00 C 1494 3497 M
-51 70 V
2257 1464 M
stroke
LT0
0.00 0.00 0.00 C 2257 1464 M
-6 91 V
stroke
LT0
0.00 0.00 0.00 C 2251 1555 M
-6 91 V
stroke
LT0
0.00 0.00 0.00 C 2245 1646 M
-6 77 V
stroke
LT0
0.00 0.00 0.00 C 2239 1723 M
0 14 V
stroke
LT0
0.00 0.00 0.00 C 2239 1737 M
-2 92 V
stroke
LT0
0.00 0.00 0.00 C 2237 1829 M
-2 91 V
stroke
LT0
0.00 0.00 0.00 C 2235 1920 M
-4 92 V
stroke
LT0
0.00 0.00 0.00 C 2231 2012 M
-5 91 V
stroke
LT0
0.00 0.00 0.00 C 2226 2103 M
-8 91 V
stroke
LT0
0.00 0.00 0.00 C 2218 2194 M
-7 59 V
stroke
LT0
0.00 0.00 0.00 C 2211 2253 M
-1 31 V
stroke
LT0
0.00 0.00 0.00 C 2210 2284 M
-9 91 V
stroke
LT0
0.00 0.00 0.00 C 2201 2375 M
-16 89 V
stroke
LT0
0.00 0.00 0.00 C 2185 2464 M
-27 88 V
stroke
LT0
0.00 0.00 0.00 C 2158 2552 M
-40 86 V
stroke
LT0
0.00 0.00 0.00 C 2118 2638 M
-28 39 V
stroke
LT0
0.00 0.00 0.00 C 2090 2677 M
-23 45 V
stroke
LT0
0.00 0.00 0.00 C 2067 2722 M
-56 83 V
stroke
LT0
0.00 0.00 0.00 C 2011 2805 M
-67 81 V
stroke
LT0
0.00 0.00 0.00 C 1944 2886 M
-51 51 V
stroke
LT0
0.00 0.00 0.00 C 1893 2937 M
-21 29 V
stroke
LT0
0.00 0.00 0.00 C 1872 2966 M
-67 81 V
stroke
LT0
0.00 0.00 0.00 C 1805 3047 M
-73 81 V
stroke
LT0
0.00 0.00 0.00 C 1732 3128 M
-23 24 V
stroke
LT0
0.00 0.00 0.00 C 1709 3152 M
-44 57 V
stroke
LT0
0.00 0.00 0.00 C 1665 3209 M
-66 82 V
stroke
LT0
0.00 0.00 0.00 C 1599 3291 M
-68 81 V
stroke
LT0
0.00 0.00 0.00 C 1531 3372 M
-12 15 V
stroke
LT0
0.00 0.00 0.00 C 1519 3387 M
-49 67 V
stroke
LT0
0.00 0.00 0.00 C 1470 3454 M
-61 83 V
2205 1430 M
stroke
LT0
0.00 0.00 0.00 C 2205 1430 M
-7 91 V
stroke
LT0
0.00 0.00 0.00 C 2198 1521 M
-7 91 V
stroke
LT0
0.00 0.00 0.00 C 2191 1612 M
-8 91 V
stroke
LT0
0.00 0.00 0.00 C 2183 1703 M
-8 91 V
stroke
LT0
0.00 0.00 0.00 C 2175 1794 M
-10 90 V
stroke
LT0
0.00 0.00 0.00 C 2165 1884 M
-11 90 V
stroke
LT0
0.00 0.00 0.00 C 2154 1974 M
-5 38 V
stroke
LT0
0.00 0.00 0.00 C 2149 2012 M
-4 53 V
stroke
LT0
0.00 0.00 0.00 C 2145 2065 M
-9 91 V
stroke
LT0
0.00 0.00 0.00 C 2136 2156 M
-13 90 V
stroke
LT0
0.00 0.00 0.00 C 2123 2246 M
-17 89 V
stroke
LT0
0.00 0.00 0.00 C 2106 2335 M
-22 88 V
stroke
LT0
0.00 0.00 0.00 C 2084 2423 M
-28 88 V
stroke
LT0
0.00 0.00 0.00 C 2056 2511 M
-37 86 V
stroke
LT0
0.00 0.00 0.00 C 2019 2597 M
-46 85 V
stroke
LT0
0.00 0.00 0.00 C 1973 2682 M
-24 34 V
stroke
LT0
0.00 0.00 0.00 C 1949 2716 M
-27 50 V
stroke
LT0
0.00 0.00 0.00 C 1922 2766 M
-54 83 V
stroke
LT0
0.00 0.00 0.00 C 1868 2849 M
-60 83 V
stroke
LT0
0.00 0.00 0.00 C 1808 2932 M
-67 80 V
stroke
LT0
0.00 0.00 0.00 C 1741 3012 M
-59 83 V
stroke
LT0
0.00 0.00 0.00 C 1682 3095 M
-62 82 V
stroke
LT0
0.00 0.00 0.00 C 1620 3177 M
-65 82 V
stroke
LT0
0.00 0.00 0.00 C 1555 3259 M
-11 12 V
stroke
LT0
0.00 0.00 0.00 C 1544 3271 M
-48 70 V
stroke
LT0
0.00 0.00 0.00 C 1496 3341 M
-60 83 V
stroke
LT0
0.00 0.00 0.00 C 1436 3424 M
-60 82 V
2162 1398 M
stroke
LT0
0.00 0.00 0.00 C 2162 1398 M
-12 91 V
stroke
LT0
0.00 0.00 0.00 C 2150 1489 M
-11 90 V
stroke
LT0
0.00 0.00 0.00 C 2139 1579 M
-4 25 V
stroke
LT0
0.00 0.00 0.00 C 2135 1604 M
-6 65 V
stroke
LT0
0.00 0.00 0.00 C 2129 1669 M
-9 91 V
stroke
LT0
0.00 0.00 0.00 C 2120 1760 M
-10 90 V
stroke
LT0
0.00 0.00 0.00 C 2110 1850 M
-12 90 V
stroke
LT0
0.00 0.00 0.00 C 2098 1940 M
-13 90 V
stroke
LT0
0.00 0.00 0.00 C 2085 2030 M
-15 90 V
stroke
LT0
0.00 0.00 0.00 C 2070 2120 M
-17 89 V
stroke
LT0
0.00 0.00 0.00 C 2053 2209 M
-21 89 V
stroke
LT0
0.00 0.00 0.00 C 2032 2298 M
-24 88 V
stroke
LT0
0.00 0.00 0.00 C 2008 2386 M
-29 88 V
stroke
LT0
0.00 0.00 0.00 C 1979 2474 M
-34 86 V
stroke
LT0
0.00 0.00 0.00 C 1945 2560 M
-40 86 V
stroke
LT0
0.00 0.00 0.00 C 1905 2646 M
-47 85 V
stroke
LT0
0.00 0.00 0.00 C 1858 2731 M
-54 83 V
stroke
LT0
0.00 0.00 0.00 C 1804 2814 M
-22 30 V
stroke
LT0
0.00 0.00 0.00 C 1782 2844 M
-32 54 V
stroke
LT0
0.00 0.00 0.00 C 1750 2898 M
-55 82 V
stroke
LT0
0.00 0.00 0.00 C 1695 2980 M
-58 83 V
stroke
LT0
0.00 0.00 0.00 C 1637 3063 M
-62 82 V
stroke
LT0
0.00 0.00 0.00 C 1575 3145 M
-3 3 V
stroke
LT0
0.00 0.00 0.00 C 1572 3148 M
-53 80 V
stroke
LT0
0.00 0.00 0.00 C 1519 3228 M
-57 83 V
stroke
LT0
0.00 0.00 0.00 C 1462 3311 M
-59 83 V
stroke
LT0
0.00 0.00 0.00 C 1403 3394 M
-39 53 V
stroke
LT0
0.00 0.00 0.00 C 1364 3447 M
-19 29 V
2122 1367 M
stroke
LT0
0.00 0.00 0.00 C 2122 1367 M
-12 90 V
stroke
LT0
0.00 0.00 0.00 C 2110 1457 M
-12 90 V
stroke
LT0
0.00 0.00 0.00 C 2098 1547 M
-13 90 V
stroke
LT0
0.00 0.00 0.00 C 2085 1637 M
-13 90 V
stroke
LT0
0.00 0.00 0.00 C 2072 1727 M
-15 90 V
stroke
LT0
0.00 0.00 0.00 C 2057 1817 M
-15 90 V
stroke
LT0
0.00 0.00 0.00 C 2042 1907 M
-14 89 V
stroke
LT0
0.00 0.00 0.00 C 2028 1996 M
-16 90 V
stroke
LT0
0.00 0.00 0.00 C 2012 2086 M
-18 89 V
stroke
LT0
0.00 0.00 0.00 C 1994 2175 M
-22 89 V
stroke
LT0
0.00 0.00 0.00 C 1972 2264 M
-25 88 V
stroke
LT0
0.00 0.00 0.00 C 1947 2352 M
-30 87 V
stroke
LT0
0.00 0.00 0.00 C 1917 2439 M
-35 86 V
stroke
LT0
0.00 0.00 0.00 C 1882 2525 M
-41 86 V
stroke
LT0
0.00 0.00 0.00 C 1841 2611 M
-41 85 V
stroke
LT0
0.00 0.00 0.00 C 1800 2696 M
-45 85 V
stroke
LT0
0.00 0.00 0.00 C 1755 2781 M
-51 84 V
stroke
LT0
0.00 0.00 0.00 C 1704 2865 M
-54 83 V
stroke
LT0
0.00 0.00 0.00 C 1650 2948 M
-46 65 V
stroke
LT0
0.00 0.00 0.00 C 1604 3013 M
-11 18 V
stroke
LT0
0.00 0.00 0.00 C 1593 3031 M
-53 83 V
stroke
LT0
0.00 0.00 0.00 C 1540 3114 M
-55 84 V
stroke
LT0
0.00 0.00 0.00 C 1485 3198 M
-57 83 V
stroke
LT0
0.00 0.00 0.00 C 1428 3281 M
-39 55 V
stroke
LT0
0.00 0.00 0.00 C 1389 3336 M
-18 28 V
stroke
LT0
0.00 0.00 0.00 C 1371 3364 M
-52 83 V
2082 1336 M
stroke
LT0
0.00 0.00 0.00 C 2082 1336 M
-12 90 V
stroke
LT0
0.00 0.00 0.00 C 2070 1426 M
-13 90 V
stroke
LT0
0.00 0.00 0.00 C 2057 1516 M
-13 90 V
stroke
LT0
0.00 0.00 0.00 C 2044 1606 M
-14 90 V
stroke
LT0
0.00 0.00 0.00 C 2030 1696 M
-14 89 V
stroke
LT0
0.00 0.00 0.00 C 2016 1785 M
-16 90 V
stroke
LT0
0.00 0.00 0.00 C 2000 1875 M
-17 89 V
stroke
LT0
0.00 0.00 0.00 C 1983 1964 M
-19 89 V
stroke
LT0
0.00 0.00 0.00 C 1964 2053 M
-21 89 V
stroke
LT0
0.00 0.00 0.00 C 1943 2142 M
-23 88 V
stroke
LT0
0.00 0.00 0.00 C 1920 2230 M
-26 88 V
stroke
LT0
0.00 0.00 0.00 C 1894 2318 M
-29 88 V
stroke
LT0
0.00 0.00 0.00 C 1865 2406 M
-32 87 V
stroke
LT0
0.00 0.00 0.00 C 1833 2493 M
-37 86 V
stroke
LT0
0.00 0.00 0.00 C 1796 2579 M
-42 85 V
stroke
LT0
0.00 0.00 0.00 C 1754 2664 M
-45 85 V
stroke
LT0
0.00 0.00 0.00 C 1709 2749 M
-51 84 V
stroke
LT0
0.00 0.00 0.00 C 1658 2833 M
-18 28 V
stroke
LT0
0.00 0.00 0.00 C 1640 2861 M
-31 56 V
stroke
LT0
0.00 0.00 0.00 C 1609 2917 M
-51 83 V
stroke
LT0
0.00 0.00 0.00 C 1558 3000 M
-52 84 V
stroke
LT0
0.00 0.00 0.00 C 1506 3084 M
-55 83 V
stroke
LT0
0.00 0.00 0.00 C 1451 3167 M
-37 53 V
stroke
LT0
0.00 0.00 0.00 C 1414 3220 M
-18 31 V
stroke
LT0
0.00 0.00 0.00 C 1396 3251 M
-51 83 V
stroke
LT0
0.00 0.00 0.00 C 1345 3334 M
-53 84 V
2049 1305 M
stroke
LT0
0.00 0.00 0.00 C 2049 1305 M
-15 90 V
stroke
LT0
0.00 0.00 0.00 C 2034 1395 M
-16 90 V
stroke
LT0
0.00 0.00 0.00 C 2018 1485 M
-10 59 V
stroke
LT0
0.00 0.00 0.00 C 2008 1544 M
-5 30 V
stroke
LT0
0.00 0.00 0.00 C 2003 1574 M
-14 90 V
stroke
LT0
0.00 0.00 0.00 C 1989 1664 M
-15 90 V
stroke
LT0
0.00 0.00 0.00 C 1974 1754 M
-16 89 V
stroke
LT0
0.00 0.00 0.00 C 1958 1843 M
-18 89 V
stroke
LT0
0.00 0.00 0.00 C 1940 1932 M
-19 89 V
stroke
LT0
0.00 0.00 0.00 C 1921 2021 M
-21 89 V
stroke
LT0
0.00 0.00 0.00 C 1900 2110 M
-24 88 V
stroke
LT0
0.00 0.00 0.00 C 1876 2198 M
-26 88 V
stroke
LT0
0.00 0.00 0.00 C 1850 2286 M
-30 88 V
stroke
LT0
0.00 0.00 0.00 C 1820 2374 M
-33 86 V
stroke
LT0
0.00 0.00 0.00 C 1787 2460 M
-37 86 V
stroke
LT0
0.00 0.00 0.00 C 1750 2546 M
-42 86 V
stroke
LT0
0.00 0.00 0.00 C 1708 2632 M
-21 40 V
stroke
LT0
0.00 0.00 0.00 C 1687 2672 M
-22 45 V
stroke
LT0
0.00 0.00 0.00 C 1665 2717 M
-44 85 V
stroke
LT0
0.00 0.00 0.00 C 1621 2802 M
-47 85 V
stroke
LT0
0.00 0.00 0.00 C 1574 2887 M
-50 83 V
stroke
LT0
0.00 0.00 0.00 C 1524 2970 M
-53 83 V
stroke
LT0
0.00 0.00 0.00 C 1471 3053 M
-29 45 V
stroke
LT0
0.00 0.00 0.00 C 1442 3098 M
-23 39 V
stroke
LT0
0.00 0.00 0.00 C 1419 3137 M
-50 84 V
stroke
LT0
0.00 0.00 0.00 C 1369 3221 M
-51 84 V
stroke
LT0
0.00 0.00 0.00 C 1318 3305 M
-52 84 V
2017 1275 M
stroke
LT0
0.00 0.00 0.00 C 2017 1275 M
-15 90 V
stroke
LT0
0.00 0.00 0.00 C 2002 1365 M
-16 90 V
stroke
LT0
0.00 0.00 0.00 C 1986 1455 M
-16 89 V
stroke
LT0
0.00 0.00 0.00 C 1970 1544 M
-17 89 V
stroke
LT0
0.00 0.00 0.00 C 1953 1633 M
-17 90 V
stroke
LT0
0.00 0.00 0.00 C 1936 1723 M
-19 89 V
stroke
LT0
0.00 0.00 0.00 C 1917 1812 M
-19 89 V
stroke
LT0
0.00 0.00 0.00 C 1898 1901 M
-7 27 V
stroke
LT0
0.00 0.00 0.00 C 1891 1928 M
-13 62 V
stroke
LT0
0.00 0.00 0.00 C 1878 1990 M
-22 88 V
stroke
LT0
0.00 0.00 0.00 C 1856 2078 M
-24 88 V
stroke
LT0
0.00 0.00 0.00 C 1832 2166 M
-27 88 V
stroke
LT0
0.00 0.00 0.00 C 1805 2254 M
-30 87 V
stroke
LT0
0.00 0.00 0.00 C 1775 2341 M
-6 16 V
stroke
LT0
0.00 0.00 0.00 C 1769 2357 M
-26 71 V
stroke
LT0
0.00 0.00 0.00 C 1743 2428 M
-34 87 V
stroke
LT0
0.00 0.00 0.00 C 1709 2515 M
-38 86 V
stroke
LT0
0.00 0.00 0.00 C 1671 2601 M
-41 86 V
stroke
LT0
0.00 0.00 0.00 C 1630 2687 M
-44 85 V
stroke
LT0
0.00 0.00 0.00 C 1586 2772 M
-47 84 V
stroke
LT0
0.00 0.00 0.00 C 1539 2856 M
-50 83 V
stroke
LT0
0.00 0.00 0.00 C 1489 2939 M
-16 27 V
stroke
LT0
0.00 0.00 0.00 C 1473 2966 M
-32 58 V
stroke
LT0
0.00 0.00 0.00 C 1441 3024 M
-49 84 V
stroke
LT0
0.00 0.00 0.00 C 1392 3108 M
-49 84 V
stroke
LT0
0.00 0.00 0.00 C 1343 3192 M
-52 84 V
stroke
LT0
0.00 0.00 0.00 C 1291 3276 M
-52 84 V
1986 1245 M
stroke
LT0
0.00 0.00 0.00 C 1986 1245 M
-16 90 V
stroke
LT0
0.00 0.00 0.00 C 1970 1335 M
-15 90 V
stroke
LT0
0.00 0.00 0.00 C 1955 1425 M
-17 89 V
stroke
LT0
0.00 0.00 0.00 C 1938 1514 M
-17 89 V
stroke
LT0
0.00 0.00 0.00 C 1921 1603 M
-18 90 V
stroke
LT0
0.00 0.00 0.00 C 1903 1693 M
-18 89 V
stroke
LT0
0.00 0.00 0.00 C 1885 1782 M
-20 88 V
stroke
LT0
0.00 0.00 0.00 C 1865 1870 M
-22 89 V
stroke
LT0
0.00 0.00 0.00 C 1843 1959 M
-22 89 V
stroke
LT0
0.00 0.00 0.00 C 1821 2048 M
-25 88 V
stroke
LT0
0.00 0.00 0.00 C 1796 2136 M
-27 87 V
stroke
LT0
0.00 0.00 0.00 C 1769 2223 M
-29 88 V
stroke
LT0
0.00 0.00 0.00 C 1740 2311 M
-31 87 V
stroke
LT0
0.00 0.00 0.00 C 1709 2398 M
-35 86 V
stroke
LT0
0.00 0.00 0.00 C 1674 2484 M
-38 86 V
stroke
LT0
0.00 0.00 0.00 C 1636 2570 M
-41 86 V
stroke
LT0
0.00 0.00 0.00 C 1595 2656 M
-44 85 V
stroke
LT0
0.00 0.00 0.00 C 1551 2741 M
-44 80 V
stroke
LT0
0.00 0.00 0.00 C 1507 2821 M
-2 5 V
stroke
LT0
0.00 0.00 0.00 C 1505 2826 M
-45 85 V
stroke
LT0
0.00 0.00 0.00 C 1460 2911 M
-46 83 V
stroke
LT0
0.00 0.00 0.00 C 1414 2994 M
-49 85 V
stroke
LT0
0.00 0.00 0.00 C 1365 3079 M
-49 84 V
stroke
LT0
0.00 0.00 0.00 C 1316 3163 M
-51 84 V
stroke
LT0
0.00 0.00 0.00 C 1265 3247 M
-5 7 V
stroke
LT0
0.00 0.00 0.00 C 1260 3254 M
-43 77 V
1955 1216 M
stroke
LT0
0.00 0.00 0.00 C 1955 1216 M
-5 21 V
stroke
LT0
0.00 0.00 0.00 C 1950 1237 M
-11 68 V
stroke
LT0
0.00 0.00 0.00 C 1939 1305 M
-16 90 V
stroke
LT0
0.00 0.00 0.00 C 1923 1395 M
-17 89 V
stroke
LT0
0.00 0.00 0.00 C 1906 1484 M
-17 89 V
stroke
LT0
0.00 0.00 0.00 C 1889 1573 M
-18 89 V
stroke
LT0
0.00 0.00 0.00 C 1871 1662 M
-19 89 V
stroke
LT0
0.00 0.00 0.00 C 1852 1751 M
-20 89 V
stroke
LT0
0.00 0.00 0.00 C 1832 1840 M
-22 89 V
stroke
LT0
0.00 0.00 0.00 C 1810 1929 M
-23 88 V
stroke
LT0
0.00 0.00 0.00 C 1787 2017 M
-25 88 V
stroke
LT0
0.00 0.00 0.00 C 1762 2105 M
-27 88 V
stroke
LT0
0.00 0.00 0.00 C 1735 2193 M
-29 87 V
stroke
LT0
0.00 0.00 0.00 C 1706 2280 M
-32 87 V
stroke
LT0
0.00 0.00 0.00 C 1674 2367 M
-35 87 V
stroke
LT0
0.00 0.00 0.00 C 1639 2454 M
-38 86 V
stroke
LT0
0.00 0.00 0.00 C 1601 2540 M
-41 85 V
stroke
LT0
0.00 0.00 0.00 C 1560 2625 M
-11 22 V
stroke
LT0
0.00 0.00 0.00 C 1549 2647 M
-30 64 V
stroke
LT0
0.00 0.00 0.00 C 1519 2711 M
-42 85 V
stroke
LT0
0.00 0.00 0.00 C 1477 2796 M
-44 85 V
stroke
LT0
0.00 0.00 0.00 C 1433 2881 M
-47 84 V
stroke
LT0
0.00 0.00 0.00 C 1386 2965 M
-48 84 V
stroke
LT0
0.00 0.00 0.00 C 1338 3049 M
-49 84 V
stroke
LT0
0.00 0.00 0.00 C 1289 3133 M
-3 5 V
stroke
LT0
0.00 0.00 0.00 C 1286 3138 M
-44 80 V
stroke
LT0
0.00 0.00 0.00 C 1242 3218 M
-47 85 V
1929 1187 M
stroke
LT0
0.00 0.00 0.00 C 1929 1187 M
-17 89 V
stroke
LT0
0.00 0.00 0.00 C 1912 1276 M
-18 89 V
stroke
LT0
0.00 0.00 0.00 C 1894 1365 M
-18 89 V
stroke
LT0
0.00 0.00 0.00 C 1876 1454 M
-18 85 V
stroke
LT0
0.00 0.00 0.00 C 1858 1539 M
-1 4 V
stroke
LT0
0.00 0.00 0.00 C 1857 1543 M
-19 89 V
stroke
LT0
0.00 0.00 0.00 C 1838 1632 M
-19 89 V
stroke
LT0
0.00 0.00 0.00 C 1819 1721 M
-20 89 V
stroke
LT0
0.00 0.00 0.00 C 1799 1810 M
-22 89 V
stroke
LT0
0.00 0.00 0.00 C 1777 1899 M
-23 88 V
stroke
LT0
0.00 0.00 0.00 C 1754 1987 M
-26 88 V
stroke
LT0
0.00 0.00 0.00 C 1728 2075 M
-27 88 V
stroke
LT0
0.00 0.00 0.00 C 1701 2163 M
-29 87 V
stroke
LT0
0.00 0.00 0.00 C 1672 2250 M
-33 87 V
stroke
LT0
0.00 0.00 0.00 C 1639 2337 M
-30 75 V
stroke
LT0
0.00 0.00 0.00 C 1609 2412 M
-4 11 V
stroke
LT0
0.00 0.00 0.00 C 1605 2423 M
-36 87 V
stroke
LT0
0.00 0.00 0.00 C 1569 2510 M
-37 86 V
stroke
LT0
0.00 0.00 0.00 C 1532 2596 M
-40 86 V
stroke
LT0
0.00 0.00 0.00 C 1492 2682 M
-42 85 V
stroke
LT0
0.00 0.00 0.00 C 1450 2767 M
-45 85 V
stroke
LT0
0.00 0.00 0.00 C 1405 2852 M
-46 84 V
stroke
LT0
0.00 0.00 0.00 C 1359 2936 M
-45 80 V
stroke
LT0
0.00 0.00 0.00 C 1314 3016 M
-2 4 V
stroke
LT0
0.00 0.00 0.00 C 1312 3020 M
-45 85 V
stroke
LT0
0.00 0.00 0.00 C 1267 3105 M
-46 85 V
stroke
LT0
0.00 0.00 0.00 C 1221 3190 M
-47 84 V
1904 1158 M
stroke
LT0
0.00 0.00 0.00 C 1904 1158 M
-18 89 V
stroke
LT0
0.00 0.00 0.00 C 1886 1247 M
-18 89 V
stroke
LT0
0.00 0.00 0.00 C 1868 1336 M
-18 89 V
stroke
LT0
0.00 0.00 0.00 C 1850 1425 M
-19 89 V
stroke
LT0
0.00 0.00 0.00 C 1831 1514 M
-20 89 V
stroke
LT0
0.00 0.00 0.00 C 1811 1603 M
-21 89 V
stroke
LT0
0.00 0.00 0.00 C 1790 1692 M
-21 88 V
stroke
LT0
0.00 0.00 0.00 C 1769 1780 M
-23 89 V
stroke
LT0
0.00 0.00 0.00 C 1746 1869 M
-24 88 V
stroke
LT0
0.00 0.00 0.00 C 1722 1957 M
-25 88 V
stroke
LT0
0.00 0.00 0.00 C 1697 2045 M
-27 88 V
stroke
LT0
0.00 0.00 0.00 C 1670 2133 M
-29 87 V
stroke
LT0
0.00 0.00 0.00 C 1641 2220 M
-31 87 V
stroke
LT0
0.00 0.00 0.00 C 1610 2307 M
-33 87 V
stroke
LT0
0.00 0.00 0.00 C 1577 2394 M
-35 87 V
stroke
LT0
0.00 0.00 0.00 C 1542 2481 M
-38 86 V
stroke
LT0
0.00 0.00 0.00 C 1504 2567 M
-40 85 V
stroke
LT0
0.00 0.00 0.00 C 1464 2652 M
-42 86 V
stroke
LT0
0.00 0.00 0.00 C 1422 2738 M
-44 85 V
stroke
LT0
0.00 0.00 0.00 C 1378 2823 M
-34 62 V
stroke
LT0
0.00 0.00 0.00 C 1344 2885 M
-11 22 V
stroke
LT0
0.00 0.00 0.00 C 1333 2907 M
-43 85 V
stroke
LT0
0.00 0.00 0.00 C 1290 2992 M
-45 85 V
stroke
LT0
0.00 0.00 0.00 C 1245 3077 M
-46 84 V
stroke
LT0
0.00 0.00 0.00 C 1199 3161 M
-47 85 V
1879 1128 M
stroke
LT0
0.00 0.00 0.00 C 1879 1128 M
-18 90 V
stroke
LT0
0.00 0.00 0.00 C 1861 1218 M
-18 89 V
stroke
LT0
0.00 0.00 0.00 C 1843 1307 M
-19 89 V
stroke
LT0
0.00 0.00 0.00 C 1824 1396 M
-19 89 V
stroke
LT0
0.00 0.00 0.00 C 1805 1485 M
-20 89 V
stroke
LT0
0.00 0.00 0.00 C 1785 1574 M
-21 88 V
stroke
LT0
0.00 0.00 0.00 C 1764 1662 M
-22 89 V
stroke
LT0
0.00 0.00 0.00 C 1742 1751 M
-22 88 V
stroke
LT0
0.00 0.00 0.00 C 1720 1839 M
-24 89 V
stroke
LT0
0.00 0.00 0.00 C 1696 1928 M
-26 88 V
stroke
LT0
0.00 0.00 0.00 C 1670 2016 M
-27 87 V
stroke
LT0
0.00 0.00 0.00 C 1643 2103 M
-29 88 V
stroke
LT0
0.00 0.00 0.00 C 1614 2191 M
-31 87 V
stroke
LT0
0.00 0.00 0.00 C 1583 2278 M
-33 87 V
stroke
LT0
0.00 0.00 0.00 C 1550 2365 M
-35 86 V
stroke
LT0
0.00 0.00 0.00 C 1515 2451 M
-38 86 V
stroke
LT0
0.00 0.00 0.00 C 1477 2537 M
-40 86 V
stroke
LT0
0.00 0.00 0.00 C 1437 2623 M
-42 85 V
stroke
LT0
0.00 0.00 0.00 C 1395 2708 M
-17 33 V
stroke
LT0
0.00 0.00 0.00 C 1378 2741 M
-25 53 V
stroke
LT0
0.00 0.00 0.00 C 1353 2794 M
-42 85 V
stroke
LT0
0.00 0.00 0.00 C 1311 2879 M
-43 84 V
stroke
LT0
0.00 0.00 0.00 C 1268 2963 M
-45 85 V
stroke
LT0
0.00 0.00 0.00 C 1223 3048 M
-46 85 V
stroke
LT0
0.00 0.00 0.00 C 1177 3133 M
-47 84 V
1853 1099 M
stroke
LT0
0.00 0.00 0.00 C 1853 1099 M
-18 90 V
stroke
LT0
0.00 0.00 0.00 C 1835 1189 M
-18 89 V
stroke
LT0
0.00 0.00 0.00 C 1817 1278 M
-19 89 V
stroke
LT0
0.00 0.00 0.00 C 1798 1367 M
-19 89 V
stroke
LT0
0.00 0.00 0.00 C 1779 1456 M
-20 89 V
stroke
LT0
0.00 0.00 0.00 C 1759 1545 M
-21 88 V
stroke
LT0
0.00 0.00 0.00 C 1738 1633 M
-22 89 V
stroke
LT0
0.00 0.00 0.00 C 1716 1722 M
-23 88 V
stroke
LT0
0.00 0.00 0.00 C 1693 1810 M
-24 88 V
stroke
LT0
0.00 0.00 0.00 C 1669 1898 M
-26 88 V
stroke
LT0
0.00 0.00 0.00 C 1643 1986 M
-27 88 V
stroke
LT0
0.00 0.00 0.00 C 1616 2074 M
-29 88 V
stroke
LT0
0.00 0.00 0.00 C 1587 2162 M
-31 87 V
stroke
LT0
0.00 0.00 0.00 C 1556 2249 M
-33 86 V
stroke
LT0
0.00 0.00 0.00 C 1523 2335 M
-36 87 V
stroke
LT0
0.00 0.00 0.00 C 1487 2422 M
-38 86 V
stroke
LT0
0.00 0.00 0.00 C 1449 2508 M
-31 68 V
stroke
LT0
0.00 0.00 0.00 C 1418 2576 M
-8 18 V
stroke
LT0
0.00 0.00 0.00 C 1410 2594 M
-39 86 V
stroke
LT0
0.00 0.00 0.00 C 1371 2680 M
-40 85 V
stroke
LT0
0.00 0.00 0.00 C 1331 2765 M
-42 85 V
stroke
LT0
0.00 0.00 0.00 C 1289 2850 M
-43 85 V
stroke
LT0
0.00 0.00 0.00 C 1246 2935 M
-45 85 V
stroke
LT0
0.00 0.00 0.00 C 1201 3020 M
-46 84 V
stroke
LT0
0.00 0.00 0.00 C 1155 3104 M
-23 43 V
1830 1071 M
stroke
LT0
0.00 0.00 0.00 C 1830 1071 M
-19 89 V
stroke
LT0
0.00 0.00 0.00 C 1811 1160 M
-11 49 V
stroke
LT0
0.00 0.00 0.00 C 1800 1209 M
-9 40 V
stroke
LT0
0.00 0.00 0.00 C 1791 1249 M
-18 89 V
stroke
LT0
0.00 0.00 0.00 C 1773 1338 M
-20 89 V
stroke
LT0
0.00 0.00 0.00 C 1753 1427 M
-20 89 V
stroke
LT0
0.00 0.00 0.00 C 1733 1516 M
-21 88 V
stroke
LT0
0.00 0.00 0.00 C 1712 1604 M
-22 89 V
stroke
LT0
0.00 0.00 0.00 C 1690 1693 M
-24 88 V
stroke
LT0
0.00 0.00 0.00 C 1666 1781 M
-24 88 V
stroke
LT0
0.00 0.00 0.00 C 1642 1869 M
-26 88 V
stroke
LT0
0.00 0.00 0.00 C 1616 1957 M
-27 88 V
stroke
LT0
0.00 0.00 0.00 C 1589 2045 M
-29 87 V
stroke
LT0
0.00 0.00 0.00 C 1560 2132 M
-32 87 V
stroke
LT0
0.00 0.00 0.00 C 1528 2219 M
-33 87 V
stroke
LT0
0.00 0.00 0.00 C 1495 2306 M
-25 61 V
stroke
LT0
0.00 0.00 0.00 C 1470 2367 M
-10 26 V
stroke
LT0
0.00 0.00 0.00 C 1460 2393 M
-35 86 V
stroke
LT0
0.00 0.00 0.00 C 1425 2479 M
-37 86 V
stroke
LT0
0.00 0.00 0.00 C 1388 2565 M
-39 86 V
stroke
LT0
0.00 0.00 0.00 C 1349 2651 M
-40 86 V
stroke
LT0
0.00 0.00 0.00 C 1309 2737 M
-42 85 V
stroke
LT0
0.00 0.00 0.00 C 1267 2822 M
-43 85 V
stroke
LT0
0.00 0.00 0.00 C 1224 2907 M
-45 84 V
stroke
LT0
0.00 0.00 0.00 C 1179 2991 M
-21 40 V
1809 1042 M
stroke
LT0
0.00 0.00 0.00 C 1809 1042 M
-19 89 V
stroke
LT0
0.00 0.00 0.00 C 1790 1131 M
-20 89 V
stroke
LT0
0.00 0.00 0.00 C 1770 1220 M
-20 89 V
stroke
LT0
0.00 0.00 0.00 C 1750 1309 M
-21 89 V
stroke
LT0
0.00 0.00 0.00 C 1729 1398 M
-21 89 V
stroke
LT0
0.00 0.00 0.00 C 1708 1487 M
-22 88 V
stroke
LT0
0.00 0.00 0.00 C 1686 1575 M
-23 89 V
stroke
LT0
0.00 0.00 0.00 C 1663 1664 M
stroke
LT0
0.00 0.00 0.00 C 1663 1664 M
-23 88 V
stroke
LT0
0.00 0.00 0.00 C 1640 1752 M
-25 88 V
stroke
LT0
0.00 0.00 0.00 C 1615 1840 M
-26 88 V
stroke
LT0
0.00 0.00 0.00 C 1589 1928 M
-14 46 V
stroke
LT0
0.00 0.00 0.00 C 1575 1974 M
-13 42 V
stroke
LT0
0.00 0.00 0.00 C 1562 2016 M
-29 87 V
stroke
LT0
0.00 0.00 0.00 C 1533 2103 M
-30 87 V
stroke
LT0
0.00 0.00 0.00 C 1503 2190 M
-31 87 V
stroke
LT0
0.00 0.00 0.00 C 1472 2277 M
-34 87 V
stroke
LT0
0.00 0.00 0.00 C 1438 2364 M
-35 86 V
stroke
LT0
0.00 0.00 0.00 C 1403 2450 M
-37 87 V
stroke
LT0
0.00 0.00 0.00 C 1366 2537 M
-39 85 V
stroke
LT0
0.00 0.00 0.00 C 1327 2622 M
-40 86 V
stroke
LT0
0.00 0.00 0.00 C 1287 2708 M
-42 86 V
stroke
LT0
0.00 0.00 0.00 C 1245 2794 M
-43 85 V
stroke
LT0
0.00 0.00 0.00 C 1202 2879 M
-17 31 V
1788 1014 M
stroke
LT0
0.00 0.00 0.00 C 1788 1014 M
-19 89 V
stroke
LT0
0.00 0.00 0.00 C 1769 1103 M
-20 89 V
stroke
LT0
0.00 0.00 0.00 C 1749 1192 M
-20 89 V
stroke
LT0
0.00 0.00 0.00 C 1729 1281 M
-21 89 V
stroke
LT0
0.00 0.00 0.00 C 1708 1370 M
-21 88 V
stroke
LT0
0.00 0.00 0.00 C 1687 1458 M
-23 89 V
stroke
LT0
0.00 0.00 0.00 C 1664 1547 M
-22 88 V
stroke
LT0
0.00 0.00 0.00 C 1642 1635 M
-24 88 V
stroke
LT0
0.00 0.00 0.00 C 1618 1723 M
-25 88 V
stroke
LT0
0.00 0.00 0.00 C 1593 1811 M
-26 88 V
stroke
LT0
0.00 0.00 0.00 C 1567 1899 M
-27 88 V
stroke
LT0
0.00 0.00 0.00 C 1540 1987 M
-29 88 V
stroke
LT0
0.00 0.00 0.00 C 1511 2075 M
-30 87 V
stroke
LT0
0.00 0.00 0.00 C 1481 2162 M
-32 87 V
stroke
LT0
0.00 0.00 0.00 C 1449 2249 M
-33 87 V
stroke
LT0
0.00 0.00 0.00 C 1416 2336 M
-35 86 V
stroke
LT0
0.00 0.00 0.00 C 1381 2422 M
-37 86 V
stroke
LT0
0.00 0.00 0.00 C 1344 2508 M
-39 86 V
stroke
LT0
0.00 0.00 0.00 C 1305 2594 M
-40 86 V
stroke
LT0
0.00 0.00 0.00 C 1265 2680 M
-42 85 V
stroke
LT0
0.00 0.00 0.00 C 1223 2765 M
-8 15 V
1767 986 M
stroke
LT0
0.00 0.00 0.00 C 1767 986 M
-19 89 V
stroke
LT0
0.00 0.00 0.00 C 1748 1075 M
-20 89 V
stroke
LT0
0.00 0.00 0.00 C 1728 1164 M
-21 88 V
stroke
LT0
0.00 0.00 0.00 C 1707 1252 M
-20 89 V
stroke
LT0
0.00 0.00 0.00 C 1687 1341 M
-22 89 V
stroke
LT0
0.00 0.00 0.00 C 1665 1430 M
-22 88 V
stroke
LT0
0.00 0.00 0.00 C 1643 1518 M
-23 89 V
stroke
LT0
0.00 0.00 0.00 C 1620 1607 M
-24 88 V
stroke
LT0
0.00 0.00 0.00 C 1596 1695 M
-25 88 V
stroke
LT0
0.00 0.00 0.00 C 1571 1783 M
-26 88 V
stroke
LT0
0.00 0.00 0.00 C 1545 1871 M
-27 88 V
stroke
LT0
0.00 0.00 0.00 C 1518 1959 M
-29 87 V
stroke
LT0
0.00 0.00 0.00 C 1489 2046 M
-30 87 V
stroke
LT0
0.00 0.00 0.00 C 1459 2133 M
-32 87 V
stroke
LT0
0.00 0.00 0.00 C 1427 2220 M
-33 87 V
stroke
LT0
0.00 0.00 0.00 C 1394 2307 M
-36 86 V
stroke
LT0
0.00 0.00 0.00 C 1358 2393 M
-37 87 V
stroke
LT0
0.00 0.00 0.00 C 1321 2480 M
-38 85 V
stroke
LT0
0.00 0.00 0.00 C 1283 2565 M
-35 73 V
1746 958 M
stroke
LT0
0.00 0.00 0.00 C 1746 958 M
-19 88 V
stroke
LT0
0.00 0.00 0.00 C 1727 1046 M
-20 89 V
stroke
LT0
0.00 0.00 0.00 C 1707 1135 M
-21 89 V
stroke
LT0
0.00 0.00 0.00 C 1686 1224 M
-21 89 V
stroke
LT0
0.00 0.00 0.00 C 1665 1313 M
-21 88 V
stroke
LT0
0.00 0.00 0.00 C 1644 1401 M
-22 89 V
stroke
LT0
0.00 0.00 0.00 C 1622 1490 M
-23 88 V
stroke
LT0
0.00 0.00 0.00 C 1599 1578 M
-24 89 V
stroke
LT0
0.00 0.00 0.00 C 1575 1667 M
-25 88 V
stroke
LT0
0.00 0.00 0.00 C 1550 1755 M
-26 87 V
stroke
LT0
0.00 0.00 0.00 C 1524 1842 M
-28 88 V
stroke
LT0
0.00 0.00 0.00 C 1496 1930 M
-29 88 V
stroke
LT0
0.00 0.00 0.00 C 1467 2018 M
-30 87 V
stroke
LT0
0.00 0.00 0.00 C 1437 2105 M
-32 87 V
stroke
LT0
0.00 0.00 0.00 C 1405 2192 M
-34 86 V
stroke
LT0
0.00 0.00 0.00 C 1371 2278 M
-35 87 V
stroke
LT0
0.00 0.00 0.00 C 1336 2365 M
-37 86 V
stroke
LT0
0.00 0.00 0.00 C 1299 2451 M
-12 28 V
1725 929 M
stroke
LT0
0.00 0.00 0.00 C 1725 929 M
-19 89 V
stroke
LT0
0.00 0.00 0.00 C 1706 1018 M
-20 89 V
stroke
LT0
0.00 0.00 0.00 C 1686 1107 M
-21 89 V
stroke
LT0
0.00 0.00 0.00 C 1665 1196 M
-21 88 V
stroke
LT0
0.00 0.00 0.00 C 1644 1284 M
-22 89 V
stroke
LT0
0.00 0.00 0.00 C 1622 1373 M
-22 89 V
stroke
LT0
0.00 0.00 0.00 C 1600 1462 M
-23 88 V
stroke
LT0
0.00 0.00 0.00 C 1577 1550 M
-24 88 V
stroke
LT0
0.00 0.00 0.00 C 1553 1638 M
-25 88 V
stroke
LT0
0.00 0.00 0.00 C 1528 1726 M
-26 88 V
stroke
LT0
0.00 0.00 0.00 C 1502 1814 M
-28 88 V
stroke
LT0
0.00 0.00 0.00 C 1474 1902 M
-29 87 V
stroke
LT0
0.00 0.00 0.00 C 1445 1989 M
-30 87 V
stroke
LT0
0.00 0.00 0.00 C 1415 2076 M
-32 87 V
stroke
LT0
0.00 0.00 0.00 C 1383 2163 M
-34 87 V
stroke
LT0
0.00 0.00 0.00 C 1349 2250 M
-15 37 V
1636 1203 M
stroke
LT0
0.00 0.00 0.00 C 1636 1203 M
-13 53 V
stroke
LT0
0.00 0.00 0.00 C 1623 1256 M
-22 89 V
stroke
LT0
0.00 0.00 0.00 C 1601 1345 M
-22 88 V
stroke
LT0
0.00 0.00 0.00 C 1579 1433 M
-24 88 V
stroke
LT0
0.00 0.00 0.00 C 1555 1521 M
-24 89 V
stroke
LT0
0.00 0.00 0.00 C 1531 1610 M
-25 88 V
stroke
LT0
0.00 0.00 0.00 C 1506 1698 M
-26 88 V
stroke
LT0
0.00 0.00 0.00 C 1480 1786 M
-28 87 V
stroke
LT0
0.00 0.00 0.00 C 1452 1873 M
-29 88 V
stroke
LT0
0.00 0.00 0.00 C 1423 1961 M
-18 51 V
stroke
LTb
4549 1228 M
1721 779 L
1119 2892 M
1721 779 L
0 1749 V
1119 4640 M
1721 2528 L
2828 448 R
1721 2528 L
stroke
LCb setrgbcolor
LTb
LCb setrgbcolor
LTb
LCb setrgbcolor
LTb
1.000 UP
grestore % colour palette end
stroke
grestore
end
showpage
  }}%
  \put(412,3653){\makebox(0,0){\strut{}f(x,y)}}%
  \put(713,1723){\makebox(0,0){\strut{}y}}%
  \put(3284,475){\makebox(0,0){\strut{}x}}%
  \put(412,3653){\makebox(0,0){\strut{}f(x,y)}}%
  \put(993,4640){\makebox(0,0)[r]{\strut{} 50}}%
  \put(993,4390){\makebox(0,0)[r]{\strut{} 40}}%
  \put(993,4140){\makebox(0,0)[r]{\strut{} 30}}%
  \put(993,3890){\makebox(0,0)[r]{\strut{} 20}}%
  \put(993,3641){\makebox(0,0)[r]{\strut{} 10}}%
  \put(993,3391){\makebox(0,0)[r]{\strut{} 0}}%
  \put(993,3141){\makebox(0,0)[r]{\strut{}-10}}%
  \put(993,2892){\makebox(0,0)[r]{\strut{}-20}}%
  \put(1002,2861){\makebox(0,0)[r]{\strut{} 2}}%
  \put(1102,2509){\makebox(0,0)[r]{\strut{} 1}}%
  \put(1203,2156){\makebox(0,0)[r]{\strut{} 0}}%
  \put(1303,1804){\makebox(0,0)[r]{\strut{}-1}}%
  \put(1403,1452){\makebox(0,0)[r]{\strut{}-2}}%
  \put(1503,1100){\makebox(0,0)[r]{\strut{}-3}}%
  \put(1603,747){\makebox(0,0)[r]{\strut{}-4}}%
  \put(4574,1076){\makebox(0,0){\strut{} 3}}%
  \put(4103,1001){\makebox(0,0){\strut{} 2}}%
  \put(3632,926){\makebox(0,0){\strut{} 1}}%
  \put(3160,851){\makebox(0,0){\strut{} 0}}%
  \put(2690,776){\makebox(0,0){\strut{}-1}}%
  \put(2218,701){\makebox(0,0){\strut{}-2}}%
  \put(1747,626){\makebox(0,0){\strut{}-3}}%
\end{picture}%
\endgroup
\endinput

  \end{center}
\end{figure}

Из графика видно, что минимум достигается близко к точке $(0, -1)$, будем исследовать функцию в окрестности этой точки.

Из того, что функция в исследуемой области не имеет особенных точек, ограничена и гладка следует, 
что выполняется условие Липшица:
$$
  \exists R \in \RR: \quad \left|\left| \nabla f(x) - \nabla f(y) \right|\right| \leqslant R || x - y ||, \forall x, y \in \RR^n,
$$
следовательно итерационный процесс градиентного спуска будет сходиться.

\subsection{Генетический алгоритм}

%FIXME:

\section{Описание алгоритма}
\subsection{Метод градиентного спуска}

\textit{Метод градиентного спуска} основывается на том, что для гладкой выпуклой функции градиент функции в точке направлен
в сторону увеличения функции (в некоторой окрестности).
Используя этот факт строится итерационный процесс приближения рассматриваемых точек области определения 
к точке минимума.

Выбирается начальное приближение минимума, далее строится последовательность точек,
в которой каждая следующая точка выбирается на антиградиенте (луче противоположном градиенту) в текущей точке:
$$ x_{k+1} = x_k - \lambda_k \nabla f(x_k), \quad \lambda_k > 0 $$

Шаг, на который ``двигается'' текущая точка за одну итерацию, выбирается следующим образом:
$$
  \lambda_k \in (0, q)\!: 
    \quad f(x_k - \lambda_k \nabla f(x_k)) = \min\limits_{0 < \lambda < q} f(x_k - \lambda \nabla f(x_k)),
$$
значение $\lambda_k$ ищется \textit{методом золотого сечения}.

Константа $q$ задаёт интервал поиска минимума на антиградиенте.

Условием остановки итерационного процесса является событие, когда следующая точка находится от предыдущей 
на расстоянии меньшим $\varepsilon$:
$$ || x_{k+1} - x_k || < \varepsilon. $$ % TODO: Module symbol.
% TODO: Критерий остановки по малой величине градиента.

%TODO!: Псевдокод алгоритма.

\subsection{Генетический алгоритм}
%FIXME:

\section{Код программы}
\subsection{Метод градиентного спуска}
\lstset{language=C++, caption=Градиентный спуск,%
label=gd-source-code, basicstyle=\footnotesize,%
numbers=left, numberstyle=\footnotesize, numbersep=5pt, frame=single, breaklines=true, breakatwhitespace=false,%
inputencoding=utf8x}
\lstinputlisting{data/gradient_descent.hpp}

\subsection{Генетический алгоритм}
%FIXME:

\section{Результаты решения}
\subsection{Метод градиентного спуска}
Результаты решения приведены в таблице \ref{gd-result-table}.

%TODO: spaces!
Начальной точкой была выбрана точка ../../input/gd_start_point.tex\!, % TODO: \!
шаг для поиска минимума методом золотого сечения был равен ../../input/gd_step.tex\!. % TODO: \!
% TODO: Рисунок с шагами.

\begin{table}[H]
\caption{Результаты работы алгоритма градиентного спуска}
\label{gd-result-table}
\begin{center}
\begin{tabular}{|c|c|c|c|c|}
\hline
Точность & Количеcтво шагов & $x$ & $f(x)$ & $f_i(x) - f_{i - 1}(x)$ \\
\hline
%TODO: Formatting.
1e-03 & 12 & (1.67515e-05,     -0.81647076) &      4.89897949 &  & (4.10337e-05, 4.744064e-05) \\
1e-04 & 12 & (5.32455e-06,     -0.81647068) &      4.89897949 & -3.053593e-10 & (1.30426e-05, 4.757985e-05) \\
1e-05 & 13 & (2.39964e-07,     -0.81649641) &      4.89897949 & -6.507568e-10 & (5.8779e-07, 3.093111e-07) \\
1e-06 & 13 & (4.23671e-07,     -0.81649656) &      4.89897949 & 1.234568e-13 & (1.03778e-06, 3.759286e-08) \\
1e-07 & 14 & (8.27773e-09,     -0.81649657) &      4.89897949 & -2.202682e-13 & (2.02762e-08, 2.704985e-08) \\
1e-08 & 15 & (2.61515e-09,     -0.81649658) &      4.89897949 & 0.000000e+00 & (6.40578e-09, 2.714819e-09) \\

\hline
\end{tabular}
\end{center}
\end{table}

\subsection{Генетический алгоритм}
%FIXME:

\begin{table}[H]
\caption{Результаты работы генетического алгоритма}
\label{gen-result-table}
\begin{center}
\begin{tabular}{|c|c|c|c|c|}
\hline
Точность & Количеcтво шагов & $x$ & $f(x)$ & $f_i(x) - f_{i - 1}(x)$ \\
\hline
%TODO: Formatting.
1e-03 & 51 & (-8.89596e-05,     -0.81641967) &      4.89897950 &  & (-0.000217887, 1.413057e-04) \\
1e-04 & 73 & (1.91949e-06,     -0.81650223) &      4.89897949 & -1.509204e-08 & (4.70178e-06, -1.037411e-05) \\
1e-05 & 89 & (1.91949e-06,     -0.81650223) &      4.89897949 & 0.000000e+00 & (4.70178e-06, -1.037411e-05) \\
1e-06 & 120 & (2.42424e-07,     -0.81649652) &      4.89897949 & -3.372858e-11 & (5.93816e-07, 1.107199e-07) \\
1e-07 & 120 & (2.42424e-07,     -0.81649652) &      4.89897949 & 0.000000e+00 & (5.93816e-07, 1.107199e-07) \\
1e-08 & 120 & (2.42424e-07,     -0.81649652) &      4.89897949 & 0.000000e+00 & (5.93816e-07, 1.107199e-07) \\

\hline
\end{tabular}
\end{center}
\end{table}

\section{Возможные дополнительные исследования}
\subsection{Метод градиентного спуска}
%FIXME:

\subsection{Генетический алгоритм}
%FIXME:

\section{Обоснование достоверности полученного результата}
\subsection{Метод градиентного спуска}
%FIXME:

\subsection{Генетический алгоритм}
%FIXME:

\end{document}
